% =====================================================================
%  capitulos/cap11.tex — Capítulo 11. Las perífrasis verbales
%  Autores: Mabel Giammatteo, Ana Marcovecchio y Andrés Kaller
%  Fuente: Capitulo_11_normalizado.docx
% =====================================================================
\capitulo{Perífrasis y locuciones verbales}{Mabel Giammatteo, Ana Marcovecchio y Andrés Kaller}

\section{Introducción}

Hasta aquí, la estructura oracional se ha presentado articulada en torno a una
única forma verbal, en la que coexisten la información léxica, aportada por la
raíz, y el contenido gramatical, expresado por la flexión. En este capítulo nos
ocupamos de los casos en que las unidades verbales son más complejas y están
constituidas por un verbo que funciona como centro léxico-semántico de la
predicación y otro que lo auxilia, manifestando la información gramatical
(tiempo, modo y aspecto) y las categorías de concordancia (persona y número).
Esas unidades que funcionan conjuntamente como núcleo del SV reciben la
denominación de \termino{perífrasis verbales} (PPVV) y, según el aporte de
significado que hace el auxiliar, se especializan en manifestar diferentes
valores relacionados con el evento: aspectuales (\forma{empezar a leer, seguir
luchando}), modales (\forma{tener que estudiar, poder descifrar}), temporales
(\forma{ir a comer, haber disfrutado}) o de diátesis pasiva (\forma{ser
premiado}). En este capítulo nos ocupamos también de los casos más complejos, en
los que se forman cadenas de auxiliares que preceden al verbo principal o
auxiliado, así como de distinguir las PPVV de las \termino{locuciones verbales}
(LLVV), que también son unidades verbales complejas, pero cuyo significado no
deriva composicionalmente del de sus constituyentes, por lo que su
interpretación no es literal (así, por ejemplo, \forma{dar fe} significa
\sig{atestiguar} y \forma{echar de menos} se interpreta como \sig{extrañar}).

\section{Concepto de perífrasis verbal}\label{sec:11-concepto}

Según \textcite[3425]{gomeztorrego1999}:

\begin{cita}
Una perífrasis verbal es la unión de dos o más verbos que constituyen un solo
núcleo del predicado. El primer verbo, llamado «auxiliar», comporta las
informaciones morfológicas de número y persona, y se conjuga en todas (o en
parte de) las formas y tiempos de la conjugación. El segundo verbo, llamado
«auxiliado» o «principal», debe aparecer en infinitivo, gerundio o participio,
es decir, es una forma no personal. Según se trate de una u otra forma, hablamos
de perífrasis de infinitivo, de gerundio o de participio. A veces contiene
partículas de enlace como \textit{a}, \textit{de}, \textit{por} y \textit{que}.
\end{cita}

Las perífrasis verbales se constituyen por un proceso de unificación de dos
verbos: el \termino{auxiliar} y el verbo \termino{principal} o
\termino{auxiliado}, que transmite el contenido léxico de la predicación. Los
auxiliares pueden ser verbos con escaso contenido semántico (por ejemplo, los
copulativos\footnote{La noción de verbo copulativo se introdujo en el
\capref{9}.} \forma{ser} y \forma{estar} y otros, como \forma{haber} y
\forma{soler}) o pueden originarse en verbos de significado léxico pleno (por
ejemplo, \forma{deber, ir, andar}) que, en su empleo como auxiliares, sufren
alguna reducción de su contenido y pasan a interpretarse con un significado
gramatical más abstracto, característico de los auxiliares. El verbo principal,
en cambio, desprovisto de flexión y en una forma no finita (infinitivo,
participio o gerundio), actúa como verbo pleno, ya que mantiene sus valores
semánticos y requisitos argumentales regulares. Así, en \exref{ex:11-fiesta}, es
\forma{comer} el que selecciona el tema ---\forma{sándwiches y bocaditos}--- y
no el auxiliar \forma{estar}; en \exref{ex:11-mozos}, ambos argumentos, tema
---\forma{los alimentos}--- y locativo ---\forma{en las mesas}---, son exigidos
por el verbo de colocación \forma{poner} y no por el auxiliar \forma{empezar}.
Además, con ambos verbos, son los auxiliados los que exigen un agente como
sujeto.

\begin{exe}
  \ex\label{ex:11-fiesta} Todos los invitados a la fiesta \textit{estaban
      comiendo} sándwiches y bocaditos.
  \ex\label{ex:11-mozos} Los mozos \textit{empezaron a poner} los alimentos en
      las mesas muy temprano.
\end{exe}

Según vemos en \exref{ex:11-mozos}, entre las dos formas verbales que integran la
perífrasis puede haber otro elemento, que puede ser una preposición, como
\forma{a} ---\forma{empezaron a correr}--- o \forma{de} ---\forma{terminaron de
decir}---, y también la conjunción \forma{que} ---\forma{tener que salir}---.

\subsection{Criterios de reconocimiento}\label{sec:11-criterios}

\subsubsection{La referencia a un evento único}\label{sec:11-evento}

Para que la unión de los dos verbos, auxiliar y auxiliado, constituya una
perífrasis, ambos deben referir a un único evento. Veamos algunos ejemplos:

\begin{exe}
  \ex\label{ex:11-juan}
    \begin{xlist}
      \ex\label{ex:11-juan-a} Juan \textit{está estudiando} en su habitación.
      \ex\label{ex:11-juan-b} Juan \textit{estudia} en su habitación.
    \end{xlist}
  \ex\label{ex:11-alumnos}
    \begin{xlist}
      \ex\label{ex:11-alumnos-a} Los alumnos \textit{deben asistir} con uniforme.
      \ex\label{ex:11-alumnos-b} Los alumnos \textit{asistirán} con uniforme.
    \end{xlist}
\end{exe}

En \exref{ex:11-juan-a}, \forma{está estudiando} equivale a la forma simple del
verbo, \forma{estudia} \exref{ex:11-juan-b}, a la que añade un matiz de
\sig{duración}; en \exref{ex:11-alumnos-a}, \forma{deben asistir} es equivalente
a \forma{asistirán} \exref{ex:11-alumnos-b}, al que agrega el valor de
\sig{obligación}. En ambos casos hay un único evento.

\subsubsection{La conmutación estructural}\label{sec:11-estructural}

Sin embargo, no siempre que existen dos verbos seguidos, uno conjugado y otro en
una forma no finita, ya sea participio, infinitivo o gerundio, se trata de una
perífrasis. Para establecer las diferencias podemos comparar las oraciones de
\exref{ex:11-problema} y \exref{ex:11-profesor}:

\begin{exe}
  \ex\label{ex:11-problema}
    \begin{xlist}
      \ex\label{ex:11-problema-a} El problema \textit{fue resuelto}.
      \ex\label{ex:11-problema-b} El problema \textit{se resolvió}. $\rightarrow$
          Un solo evento.
      \ex[*]{\label{ex:11-problema-c} El problema fue que resuelto.}
          Rechazo de la conmutación estructural.
    \end{xlist}
  \ex\label{ex:11-profesor}
    \begin{xlist}
      \ex\label{ex:11-profesor-a} El profesor \textit{decidió posponer} el
          examen.
      \ex\label{ex:11-profesor-b} El profesor \textit{decidió (tomó la decisión) y
          pospuso} el examen. $\rightarrow$ Dos eventos.
      \ex\label{ex:11-profesor-c} El profesor \textit{decidió que posponía} el
          examen. Aceptación de la conmutación estructural.
    \end{xlist}
\end{exe}

Una prueba que podemos aplicar para distinguir estos casos es la de
\termino{conmutación}\footnote{Se entiende por \textit{conmutación} la operación
que consiste en cambiar un elemento por un equivalente funcional para comprobar
si produce o no cambio de significado. Por ejemplo, \forma{Llegó triste} puede
ser reemplazado por \forma{Llegó con tristeza} o \forma{Llegó apenado}. En el
\capref{3} esta prueba se aplicó para contrastar sonidos y corroborar si se
trataba de fonemas diferentes o de alófonos del mismo fonema.}
\termino{estructural}, que consiste en cambiar la estructura sintáctica para ver
si se trata de dos formas unificadas, como en las perífrasis, o de dos formas
independientes, cada una de las cuales remite a un evento diferente. En
\exref{ex:11-problema-a} tenemos una perífrasis porque ambas formas verbales
remiten a un único evento, de modo que es posible cambiar \forma{fue resuelto}
por otra forma simple equivalente, \forma{se resolvió} \exref{ex:11-problema-b};
pero un cambio estructural, como el que aparece en \exref{ex:11-problema-c},
produce un resultado agramatical. El caso de \exref{ex:11-profesor-a}, en cambio,
no tiene una interpretación de evento único, como vemos en
\exref{ex:11-profesor-b}, y, por lo tanto, acepta el cambio por otra estructura
\exref{ex:11-profesor-c}, lo que nos permite comprobar que no se trata de una
perífrasis.

En las PPVV, auxiliar y auxiliado forman una unidad compleja que funciona como
núcleo del SV \exref{ex:11-problema-a}, como vemos en \exref{ex:11-de-pv}:

\begin{exe}
  \ex\label{ex:11-de-pv} {}[\etq{SF/O} [\etq{SN/Suj} El problema] [\etq{SV/Pred}
      [\etq{V/n} fue resuelto] [\etq{SP/CAg} por los estudiantes]]]
\end{exe}

En cambio, en el caso de \exref{ex:11-profesor-a}, cada uno de los verbos
pertenece a una oración diferente: la forma conjugada es el núcleo del predicado
verbal de la oración principal y la forma no finita constituye el verbo núcleo de
otra oración, subordinada,\footnote{El concepto de oración subordinada se
desarrolla en el \capref{13}.} tal como vemos en \exref{ex:11-de-nopv}:

\begin{exe}
  \ex\label{ex:11-de-nopv} {}[\etq{SF/O} [\etq{SN/Suj} El profesor]
      [\etq{SV/Pred} decidió\etq{V/n} [\etq{SC=SF/COD} posponer el
      examen]]]\pendiente{El original analiza aquí «resolvió aplazar»}
\end{exe}

En algunos casos de lectura ambigua, como \exref{ex:11-perdon-a}, la conmutación
estructural permite anclar una interpretación: en \exref{ex:11-perdon-b} se
interpreta \forma{volvieron a pedir} como un evento que se repite (=
\sig{pidieron de nuevo}), por lo que se trata de una perífrasis; pero en
\exref{ex:11-perdon-c}, con \forma{volver} acompañado de un complemento locativo
---\forma{volvieron al lugar}---, se disuelve la posibilidad de lectura como
perífrasis, ya que admite conmutación estructural: el infinitivo se puede
reemplazar por un SP encabezado por \forma{para} ---\forma{volvieron para pedir
perdón}---.

\begin{exe}
  \ex
    \begin{xlist}
      \ex\label{ex:11-perdon-a} \textit{Volvieron a pedir} perdón.
      \ex\label{ex:11-perdon-b} Reiteraron sus disculpas / Pidieron perdón de
          nuevo. $\rightarrow$ Un solo evento.
      \ex\label{ex:11-perdon-c} Volvieron al lugar y pidieron perdón.
          $\rightarrow$ Dos eventos.
    \end{xlist}
\end{exe}

\begin{ejercicio}\label{ej:11-reconocer}
En las siguientes oraciones, determine si se trata o no de PPVV o si hay
ambigüedad. Según se muestra en los ejemplos resueltos de abajo, fundamente sus
respuestas aplicando las pruebas de reconocimiento suministradas (referencia a
un evento único, admisión o no de conmutación estructural). En los casos
ambiguos, para las PPVV ofrezca una paráfrasis adecuada del significado.

\begin{exe}
  \ex La sentencia \textit{ha salido} del juzgado hoy.
  \ex Los deportistas \textit{vienen corriendo} más rápido.
  \ex Mi abuela no \textit{salió a pasear} en la pandemia.
  \ex \textit{Volvió a hacer} el trabajo práctico de la materia.
  \ex El domingo \textit{empezó a llover} desde temprano.
\end{exe}

\begin{resuelto}[Ejemplos resueltos]
1. Es PV porque señala un único evento y no admite conmutación estructural:
\juicio{*}{\forma{La sentencia ha y salió del juzgado hoy}}. 2. Es ambigua y
posee dos interpretaciones: i) no es PV porque indica dos eventos y admite
conmutación estructural: \forma{Los deportistas vienen (al club) y en el
trayecto apuran su marcha}; ii) es PV y no admite conmutación estructural.
Paráfrasis: \forma{Los deportistas mejoran progresivamente sus habilidades para
la carrera}.
\end{resuelto}
\end{ejercicio}

\subsubsection{La conmutación léxica}\label{sec:11-lexica}

Aunque no admiten conmutación estructural, las PPVV aceptan \termino{conmutación
léxica}. Esta prueba, que consiste en cambiar una palabra por otra de la misma
categoría, se aplica a las PPVV mediante la conmutación del verbo principal por
otro de significado equivalente:

\begin{exe}
  \ex
    \begin{xlist}
      \ex Ana \textit{debe comer} mejor porque está anémica.
      \ex Ana \textit{debe alimentarse} mejor porque está anémica.
    \end{xlist}
  \ex
    \begin{xlist}
      \ex El famoso actor \textit{fue premiado} por su labor en la nueva serie
          de televisión.
      \ex El famoso actor \textit{fue galardonado} por su labor en la nueva
          serie de televisión.
    \end{xlist}
  \ex
    \begin{xlist}
      \ex Todos \textit{continuaron mirando} apasionadamente la nueva temporada
          de la serie.
      \ex Todos \textit{continuaron viendo} apasionadamente la nueva temporada
          de la serie.
    \end{xlist}
\end{exe}

\begin{ejercicio}\label{ej:11-series}
Lea atentamente el texto siguiente y resuelva lo pedido en las consignas que se
presentan debajo.

% Texto periodístico protegido: se incorpora a mano desde el original.
\textoaincorporar{el fragmento adaptado de \textit{Infobae} (08/11/2020) sobre
dos tipos de series, con las formas en cursiva: \textit{nos hacen sentir; son
devoradas; puede distraer; comenzarlas a ver; permiten conectarnos; suelen
proveer}}

\begin{consignas}
  \item Determine si las formas en cursiva constituyen o no PPVV. Aporte pruebas
        que fundamenten sus afirmaciones (referencia o no a un evento único,
        conmutación léxica y conmutación estructural), según se muestra en el
        ejemplo resuelto~1.
  \item Reescriba el fragmento incorporando los sinónimos obtenidos en la
        conmutación léxica, sin alterar el significado del texto.
\end{consignas}

\begin{resuelto}
1. \forma{(Nos) hacen sentir} no es PV. Se trata de dos eventos; admite
conmutación léxica, \forma{(nos) hacen experimentar}, y conmutación estructural,
\forma{hacen que sintamos}.
\end{resuelto}
\end{ejercicio}

\section{Criterios de clasificación}\label{sec:11-clasificacion}

Las perífrasis verbales suelen clasificarse teniendo en cuenta dos criterios
complementarios: a) la forma no finita que interviene en la perífrasis y b) el
aporte semántico del auxiliar.

\begin{itemize}
  \item Según la forma no finita del verbo auxiliado,\pendiente{El original dice
        «la forma no finita del auxiliar que acompaña al verbo auxiliado»}
        tenemos:
        \begin{itemize}
          \item \termino{perífrasis de infinitivo}: \forma{volver a leer, empezar a
                escribir, poder entender}, etc.;
          \item \termino{perífrasis de gerundio}: \forma{estar jugando, seguir
                estudiando}, etc.;
          \item \termino{perífrasis de participio}: \forma{haber leído, ser
                vacunado}, etc.
        \end{itemize}
\end{itemize}

\begin{destacado}
\textbf{Para tener en cuenta.} El comportamiento del participio en cuanto a la
concordancia no es el mismo en todas las perífrasis: mientras que con los tiempos
compuestos es invariable \exref{ex:11-primos}, en la perífrasis de voz pasiva el
participio concuerda en género y número con el sujeto de la oración
\exref{ex:11-francia} y, en las perífrasis con verbos como \forma{tener} o
\forma{llevar}, concuerda con el objeto \exref{ex:11-libros}.

\begin{exe}
  \ex\label{ex:11-primos} Participio invariable: Mis primos / mis primas
      \textit{han visitado} Francia / varios países.
  \ex\label{ex:11-francia} Participio concordado con el sujeto:
    \begin{xlist}
      \ex Francia \textit{fue visitada} por mis primos.
      \ex El país de Luis XV \textit{fue visitado} por mis primos.
    \end{xlist}
  \ex\label{ex:11-libros} Participio concordado con el objeto:
    \begin{xlist}
      \ex \textit{Tengo / Llevo leídos} más de 20 libros sobre el tema.
      \ex \textit{Tengo / Llevo leídas} más de 20 novelas de este autor.
    \end{xlist}
\end{exe}
\end{destacado}

\begin{itemize}
  \item Según el aporte semántico del auxiliar, se reconocen:
        \begin{itemize}
          \item perífrasis de diátesis: \forma{ser leído};
          \item perífrasis aspectuales: \forma{empezar a leer, continuar
                escuchando};
          \item perífrasis temporales: \forma{ir a comer, haber dicho};
          \item perífrasis modales: \forma{deber escribir, tener que mirar, poder
                ser}.
        \end{itemize}
\end{itemize}

\begin{recuadro}[Para ampliar]
Con respecto a las categorías verbales (tiempo, aspecto y modo), las formas no
finitas solo manifiestan aspecto, no de forma flexional, sino inherente, es
decir, por su propia naturaleza. Como señala \textcite{hernanz1999}, el
infinitivo es neutro en relación con las otras dos formas no finitas, que marcan
valores específicos: el gerundio presenta el proceso en desarrollo, por lo
tanto, es durativo, y el participio indica el proceso como terminado; de ahí su
sentido perfectivo-resultativo. Del valor aspectual de cada forma no finita se
deriva su utilización en las PPVV:

\begin{itemize}
  \item el \termino{infinitivo}, que es neutro, es el que más combinaciones
        admite y puede manifestar valores tales como el inceptivo
        ---\forma{empezó a cantar}---, el terminativo ---\forma{dejó de
        cantar}---, el habitual ---\forma{suele cantar}--- o el de repetición
        ---\forma{volvió a cantar}---;
  \item el \termino{gerundio} se emplea en perífrasis con valor durativo, como
        la progresiva ---\forma{estoy cantando}--- y la continuativa
        ---\forma{siguió cantando}---;
  \item el \termino{participio}, que es perfectivo-resultativo, se utiliza en los
        tiempos compuestos ---\forma{había cantado}--- y en la voz pasiva
        ---\forma{fue cantado}--- \parencite[50--51]{giammatteo2009}.
\end{itemize}
\end{recuadro}

\begin{ejercicio}\label{ej:11-tato}
Considere los siguientes fragmentos y resuelva lo pedido en las consignas, según
se muestra en el ejemplo resuelto~1.

\begin{consignas}
  \item Determine si las combinaciones de palabras subrayadas son perífrasis
        verbales o no.
  \item En las que son perífrasis, haga la descripción estructural de sus
        constituyentes.
\end{consignas}

% Textos protegidos (monólogos de Tato Bores): se incorporan a mano desde el original.
\textoaincorporar{los seis fragmentos adaptados de monólogos de Tato Bores
(cultura.gob.ar, 12/01/2022), con las combinaciones subrayadas: 1. \textit{se
tuvieron que aguantar; (no) crea que le miento; voy a detallar}; 2.
\textit{iremos a comprar; van a tener que entregar; dicen ser; están corriendo};
3. \textit{estamos viviendo; estamos diciendo}; 4. \textit{estoy aburrida; dicen
defenderme}; 5. \textit{iremos a pagarle; van a tener que entregar}; 6.
\textit{estamos pagando; vamos a pagar}}

\begin{resuelto}
1. \forma{Se tuvieron que aguantar} es PV. Descripción estructural (DE): V
auxiliar \forma{tener} + \forma{que} + V auxiliado en infinitivo
\forma{aguantarse}. \forma{(No) crea que le miento}: no es PV. \forma{Voy a
detallar}: es PV. DE: V auxiliar \forma{ir} + \forma{a} + V auxiliado en
infinitivo \forma{detallar}.
\end{resuelto}
\end{ejercicio}

\section{Principales características de los auxiliares de perífrasis}\label{sec:11-auxiliares}

Los auxiliares de perífrasis verbales presentan las siguientes características:

\paragraph{Carácter delimitado de la clase.} Los verbos auxiliares que componen
las perífrasis verbales son un conjunto relativamente acotado de pocos miembros
que se repiten en la mayoría de las lenguas. Según se plantea en
\textcite[205]{giammatteo2018marcovecchio}:

\begin{cita}
[\ldots] la comparación entre lenguas no necesariamente emparentadas entre sí
nos muestra que no cualquier verbo deviene auxiliar. Así, no esperamos, ni
encontramos en las distintas lenguas del mundo, verbos con significados
específicos como, por ejemplo, \sig{comer}, \sig{leer} o \sig{cocinar}, que
hayan adoptado una función auxiliar. Por lo general, los auxiliares son verbos
de significados más generales (Heine, 1991), como \textit{do} (ingl.),
\textit{hacer} (esp.); o bien que señalan \sig{posesión}, como \textit{tener}
(esp.), \textit{have} (ingl.) y \textit{avoir} (fr.). Este tipo de verbos son
buenos candidatos para la auxiliarización.
\end{cita}

Además de los mencionados, otro grupo de verbos que frecuentemente se emplean
como auxiliares en las distintas lenguas son los de movimiento ---\forma{ir,
llevar, andar, venir, seguir}---, que pasan de señalar \sig{desplazamiento} en el
orden espacial a referir al \sig{desarrollo} del evento, ya sea en el tiempo
\parencite{cuenca1999} ---\forma{voy a comer} (valor futuro: \forma{comeré})---
o en relación con diferentes valores aspectuales ---\forma{voy / vengo / ando /
sigo / continúo estudiando}---.

\paragraph{Desemantización.} Algunos de los auxiliares, como \forma{ser, estar} y
\forma{haber}, son verbos con muy poco significado léxico inherente; otros, como
\forma{ir, venir} o \forma{seguir}, son verbos de significado pleno, que en sus
usos auxiliares se considera que han sufrido un proceso de
\termino{desemantización}. Según lo explica \textcite[21]{garciafernandez2006},
«[d]esde el punto de vista semántico, las construcciones perifrásticas suelen ser
descritas por el proceso de desemantización que experimenta el auxiliar y que
puede ilustrarse fácilmente con la oración \textit{Va a nevar}, donde la forma
\textit{va} no denota movimiento alguno». En el caso de estos auxiliares
desemantizados, es necesario destacar que en sus otros usos continúan
funcionando como verbos plenos y que pierden (parte de) su significado solo
cuando se combinan con otra forma verbal no finita, es decir que «la pérdida de
su significado viene provocada por factores colocacionales\ldots»
\parencite[124]{tsutahara2022}.

Para determinar si una secuencia de verbo conjugado más otro no finito constituye
una perífrasis verbal o si se trata de una secuencia de dos verbos, es
indispensable evaluar si el verbo conjugado mantiene su significado léxico pleno
o si sufrió algún tipo de modificación o pérdida. Por ejemplo, en
\exref{ex:11-mariano}, \forma{va} no indica \sig{desplazamiento}, sino
\sig{acción futura} y, en \exref{ex:11-sigue}, \forma{seguir} perdió el valor
espacio-temporal de \sig{ir detrás o después} para pasar a significar
\sig{continuar}.

\begin{exe}
  \ex\label{ex:11-mariano} Mariano \textit{va a publicar} su tesis.
  \ex\label{ex:11-sigue} \textit{Sigue trabajando} en una gran empresa.
\end{exe}

Sin embargo, en otros casos, la oposición entre verbos plenos y verbos
auxiliares «puede ser demasiado radical, ya que no deja lugar para los verbos que
se comportan parcialmente como auxiliares» \parencite{bosque2019}. Por ejemplo,
se considera que \forma{querer} en \exref{ex:11-salsa-a} no constituye una
perífrasis verbal porque no experimenta una interpretación distinta de su
significado de \sig{desear}; de ahí que estemos frente a un empleo de
\forma{querer} como verbo pleno y que, por lo tanto, esperaríamos que acepte la
prueba de conmutación estructural. Sin embargo, como muestra
\exref{ex:11-salsa-b}, el resultado de la conmutación estructural es una oración
agramatical.

\begin{exe}
  \ex\label{ex:11-salsa}
    \begin{xlist}
      \ex\label{ex:11-salsa-a} Yo \textit{quiero bailar} salsa.
      \ex[*]{\label{ex:11-salsa-b} Yo \textit{quiero} que (yo) \textit{baile}
          salsa.}
    \end{xlist}
\end{exe}

No obstante, lo que realmente impide aplicar la prueba de conmutación estructural
en \exref{ex:11-salsa-b} es la coincidencia de sujeto entre \forma{querer} y
\forma{bailar}: ambos verbos tienen el mismo sujeto, \forma{yo}. Cuando no hay
coincidencia de persona entre el sujeto de la forma personal y el de la forma no
finita, como se muestra en \exref{ex:11-ana}, \forma{querer} + infinitivo acepta
conmutación estructural.

\begin{exe}
  \ex\label{ex:11-ana} Yo \textit{quiero} que Ana \textit{baile} salsa.
\end{exe}

En \exref{ex:11-ana}, como se trata de dos personas gramaticales distintas para
cada verbo ---\forma{yo} para \forma{quiero} y \forma{ella} para
\forma{baile}---, la conmutación estructural es posible, por lo que queda claro
que se trata de dos eventos distintos y que no hay perífrasis. Pero ¿qué sucede
en \exref{ex:11-llover}?

\begin{exe}
  \ex\label{ex:11-llover} Parece que \textit{quiere llover}.
\end{exe}

En \exref{ex:11-llover} no le podemos asignar el deseo a ningún individuo. Por
tanto, asumimos que \forma{querer} ya no actúa como un verbo pleno que selecciona
a su sujeto, sino que ha adquirido un significado diferente al de \sig{desear}:
el de sugerir \sig{la inminencia de la lluvia}. Ambas características hacen que
en su empleo con verbos impersonales \forma{querer} se considere un auxiliar.
Como plantea \textcite[241]{ditullio2014}, «los auxiliares son semánticamente
transparentes» y, por eso, a diferencia de los verbos plenos, que imponen
restricciones a sus sujetos \exref{ex:11-escuela}, pueden construirse con verbos
impersonales y aceptar cualquier tipo de SN como sujeto
\exref{ex:11-complicado}.

\begin{exe}
  \ex\label{ex:11-escuela} Juan / *el año \textit{va} temprano a la escuela.
  \ex\label{ex:11-complicado} Juan / el año \textit{va a continuar} complicado.
\end{exe}

\paragraph{Integración con el verbo principal.} El grado de integración entre el
verbo auxiliar y el auxiliado se desprende de ciertas propiedades que actúan
sobre el conjunto, como el alcance de los adjuntos. Cuando tenemos dos verbos
independientes que no forman perífrasis, cada uno de ellos puede recibir
modificadores adjuntos \exref{ex:11-mariana1}; en las perífrasis, en cambio,
auxiliar y auxiliado deben ser modificados conjuntamente \exref{ex:11-mariana2}.

\begin{exe}
  \ex\label{ex:11-mariana1} Hoy Mariana decidió \textit{ir} mañana a la reunión en
      la facultad.
  \ex\label{ex:11-mariana2}
    \begin{xlist}
      \ex Hoy Mariana \textit{va a ir} a la facultad.
      \ex[*]{Hoy Mariana \textit{va a ir} mañana a la facultad.}
    \end{xlist}
\end{exe}

Asimismo, a diferencia de lo que sucede con la secuencia de dos verbos
\exref{ex:11-neg-a}, el auxiliar no admite ser negado en forma independiente del
verbo principal \exref{ex:11-neg-c}.\footnote{No obstante, según veremos en
§\ref{sec:11-modales}, \forma{poder} epistémico admite doble negación: en el
verbo principal y en el auxiliar.}

\begin{exe}
  \ex
    \begin{xlist}
      \ex\label{ex:11-neg-a} Mariana no decidió no \textit{ir} mañana a la
          facultad.
      \ex Mariana no \textit{va a ir} a la facultad.
      \ex[*]{\label{ex:11-neg-c} Mariana no \textit{va a} no \textit{ir} a la
          facultad.}
    \end{xlist}
\end{exe}

Otra muestra del grado de integración entre el auxiliar y el verbo principal es
que, a diferencia de lo que sucede cuando se trata de una secuencia de dos verbos
independientes \exref{ex:11-imputado}, las perífrasis no aceptan la omisión del
verbo principal ni en la interrogación \exref{ex:11-suele-b} ni en la
focalización con estructuras de relieve\footnote{El concepto de foco y los
recursos mediante los que se manifiesta se presentaron en el \capref{6},
§\ref{sec:06-capas}.} \exref{ex:11-suele-d}. No obstante, si bien en las PPVV el
verbo principal no puede elidirse, puede ser reemplazado por \forma{hacer}, tanto
en la interrogación \exref{ex:11-suele-c} como en la focalización
\exref{ex:11-suele-e}.

\begin{exe}
  \ex\label{ex:11-imputado}
    \begin{xlist}
      \ex El imputado \textit{prometió} decir la verdad.
      \ex ¿Qué \textit{prometió} el imputado? $\rightarrow$ Interrogación con
          omisión del verbo principal.
      \ex Lo que el imputado \textit{prometió} es decir la verdad. $\rightarrow$
          Focalización con omisión del verbo principal.
    \end{xlist}
  \ex\label{ex:11-suele}
    \begin{xlist}
      \ex El imputado \textit{suele mentir} con gran facilidad.
      \ex[*]{\label{ex:11-suele-b} ¿Qué \textit{suele} el imputado con gran
          facilidad?} $\rightarrow$ Interrogación sin aceptar la omisión del verbo
          principal.
      \ex\label{ex:11-suele-c} ¿Qué \textit{suele hacer} el imputado con gran
          facilidad? $\rightarrow$ Interrogación con reemplazo del verbo principal
          por el proverbo \forma{hacer}.
      \ex[*]{\label{ex:11-suele-d} Lo que el imputado \textit{suele} es mentir con
          gran facilidad.} $\rightarrow$ Focalización sin aceptar la omisión del
          verbo principal.
      \ex\label{ex:11-suele-e} Lo que el imputado \textit{suele hacer} es mentir
          con gran facilidad. $\rightarrow$ Focalización con reemplazo del verbo
          principal por el proverbo \forma{hacer}.
    \end{xlist}
\end{exe}

No obstante la cohesión que existe entre los dos verbos, es posible la
intercalación del sujeto \exref{ex:11-oferta-a} y de determinados adverbios,
\exref{ex:11-oferta-b} y \exref{ex:11-oferta-c}, entre ambos
\parencite[243]{bosque2019}.

\begin{exe}
  \ex
    \begin{xlist}
      \ex\label{ex:11-oferta-a} ¿\textit{Puede} Juan \textit{rechazar} una oferta
          tan atractiva?
      \ex\label{ex:11-oferta-b} ¿Juan \textit{puede} realmente \textit{rechazar}
          una oferta tan atractiva?
      \ex\label{ex:11-oferta-c} \textit{Terminaron} recientemente \textit{de
          elaborar} el nuevo proyecto.
    \end{xlist}
\end{exe}

\paragraph{Comportamiento respecto de los clíticos.} En las oraciones con un
núcleo verbal simple, los clíticos de COD (\forma{lo, la, los, las}) y de COI
(\forma{le, les, se}) adoptan la posición preverbal.

\begin{exe}
  \ex
    \begin{xlist}
      \ex Ana \textit{cortó} la torta. / Ana la \textit{cortó}.
      \ex Ana le \textit{dio} una porción a Juan. / Ana se la \textit{dio}.
    \end{xlist}
\end{exe}

Las perífrasis de infinitivo y de gerundio, al actuar como una unidad verbal,
aceptan que el clítico se ubique antes del verbo conjugado, pero también admiten
que se ubique después de la forma no finita, \exref{ex:11-torta} y
\exref{ex:11-examenes}. Las perífrasis de participio, en cambio, solo aceptan la
anteposición de los clíticos \exref{ex:11-premio}.

\begin{exe}
  \ex\label{ex:11-torta} Ana \textit{comenzó a cortar} la torta. / Ana la
      \textit{comenzó a cortar}. / Ana \textit{comenzó a cortarla}.
  \ex\label{ex:11-examenes} Pedro \textit{está preparando} sus exámenes. / Pedro
      los \textit{está preparando}. / Pedro \textit{está preparándolos}.
  \ex\label{ex:11-premio} El director de la compañía \textit{ha recibido} un
      premio. / El director de la compañía lo \textit{ha recibido}. /
      \juicio{*}{\forma{El director de la compañía ha recibídolo.}}
\end{exe}

A diferencia de las PPVV, las secuencias verbales no perifrásticas no tienen un
comportamiento uniforme: algunos verbos, como \forma{querer, saber, necesitar,
intentar, tratar de}, aceptan anteposición y posposición \exref{ex:11-turistas},
pero otros, como \forma{decidir, lamentar, creer, pensar}, solo admiten
pronombres enclíticos \exref{ex:11-lamento}.

\begin{exe}
  \ex\label{ex:11-turistas} Los turistas \textit{quieren conocer} la ciudad. /
      Los turistas la \textit{quieren conocer}. / Los turistas \textit{quieren
      conocerla}.
  \ex\label{ex:11-lamento} Juan \textit{lamentó comprar} ese auto. / Juan
      \textit{lamentó comprarlo}. / \juicio{*}{\forma{Juan lo lamentó comprar.}}
\end{exe}

Para \textcite{ditullio2014}, la posición de los clíticos sirve como prueba
negativa para reconocer una perífrasis, porque si una forma finita seguida de un
infinitivo no permite la anteposición, no es perífrasis \exref{ex:11-lamento};
pero si la permite, puede tanto ser \exref{ex:11-torta} como no ser
\exref{ex:11-turistas} una verdadera perífrasis.

\paragraph{Pasivización.} Del mismo modo que con los verbos simples
correspondientes, las perífrasis verbales conformadas por un verbo principal que
tenga un sujeto y un objeto directo pueden pasivizarse:

\begin{exe}
  \ex
    \begin{xlist}
      \ex Juan \textit{suele bailar} \textit{La cumparsita}.
      \ex \textit{La cumparsita} \textit{suele ser bailada} por Juan.
    \end{xlist}
\end{exe}

En cambio, en el caso de las oraciones con un verbo conjugado más un verbo no
finito que no conforman perífrasis verbal, puede suceder que, si el verbo exige
un sujeto animado, la pasiva resulte agramatical:

\begin{exe}
  \ex
    \begin{xlist}
      \ex Juan \textit{quiere bailar} \textit{La cumparsita}.
      \ex[*]{\textit{La cumparsita} \textit{quiere ser bailada} por Juan.}
    \end{xlist}
\end{exe}

\begin{ejercicio}\label{ej:11-vf}
Indique si las siguientes afirmaciones son verdaderas o falsas y justifique sus
respuestas, según se muestra en el ejemplo resuelto~1.

\begin{exe}
  \ex En la mayoría de las lenguas, los verbos auxiliares que componen
      perífrasis verbales pertenecen a un conjunto relativamente cerrado de pocos
      miembros.
  \ex Los auxiliares de perífrasis verbales son verbos desemantizados, que en
      otros contextos pueden mantener su significado y actuar como verbos
      plenos.
  \ex Las perífrasis verbales de participio permiten la posposición del clítico.
  \ex Los auxiliares de perífrasis nunca pueden construirse con verbos
      impersonales y aceptan cualquier tipo de SN como sujeto.
\end{exe}

\begin{resuelto}
1. Verdadero. La mayoría de las lenguas suele recurrir a verbos con significados
generales, como \forma{ser, tener} o \forma{hacer}, y resulta poco probable que
utilicen como auxiliares a aquellos con significados más específicos, como
\forma{estudiar, saltar} o \forma{construir}.
\end{resuelto}
\end{ejercicio}

\begin{ejercicio}\label{ej:11-caracteristicas}
Teniendo en cuenta las principales características de los auxiliares de
perífrasis verbales señaladas en §\ref{sec:11-auxiliares}, explique por qué las
combinaciones de verbos presentes en las siguientes oraciones son o no PPVV,
según se muestra en el ejemplo resuelto~1.

\begin{exe}
  \ex Juan no querrá no ir a la escuela mañana.
  \ex ¿Qué prometió Juan? (Juan prometió viajar a Córdoba).
  \ex ¿Qué suele hacer Juan todas las mañanas?
  \ex Juan pudo recientemente comprar su casa.
  \ex María necesita saberlo todo.
\end{exe}

\begin{resuelto}
1. No es PV porque no existe integración entre los verbos y cada uno de ellos
aparece negado en forma independiente.
\end{resuelto}
\end{ejercicio}

\begin{recuadro}[Para ampliar: los auxiliares de perífrasis y los procesos de gramaticalización]
En las lenguas, los verbos más frecuentes suelen sufrir un «desgaste» que va
erosionando sus significados; así, por ejemplo, el verbo \forma{habere}, que en
latín indicaba \sig{posesión}, en las lenguas romances ha perdido parcialmente
ese valor y se ha convertido en un auxiliar para los tiempos compuestos: esp.
\forma{haber}, fr. \forma{avoir}, it. \forma{avere}, rum. \forma{avea}. En estos
casos se dice que el verbo ha sufrido un proceso de
\termino{gramaticalización}, entendido como «un cambio lingüístico por el que una
voz con significado léxico adquiere valor gramatical en su evolución histórica»
\parencite[147]{bosque2019}. En el caso de \forma{haber}, en su paso del latín
al español, su significado de \sig{posesión} se ha perdido completamente.
Asimismo, también puede suceder que, como resultado de un proceso de
gramaticalización, un verbo ya de muy poco contenido, como \forma{ser}, que,
según vimos en el \capref{9}, exige un predicativo para completar su
significado, «adquiera una función aún más gramatical»
\parencite{company2003}, como cuando se convierte en el soporte de las
categorías de tiempo, modo, aspecto y concordancia (número y persona) en la voz
pasiva. En otros casos, lo que sucede es que el verbo mantiene usos con
significado pleno (\forma{Anda por toda la ciudad en moto} = \sig{se desplaza})
y otros en los que se ha desemantizado, es decir, ha perdido parte de su
significado léxico (\forma{Anda rumoreando que habrá cambios en el gobierno}).
En este último ejemplo, \forma{andar} ya no implica \sig{desplazamiento}, pero
mantiene un valor referido a que se trata de un evento continuativo-reiterativo
\parencite{yllera1999}.

Según vemos, entonces, no en todas las perífrasis el auxiliar muestra el mismo
grado de pérdida de significado léxico; por eso las características de las
perífrasis verbales que se han venido señalando dependen en gran medida del grado
de consolidación de la estructura perifrástica, condicionada, a su vez, por el
estadio de gramaticalización que haya alcanzado el auxiliar. Es decir, nos
enfrentamos a un \textit{continuum} en el que en un extremo encontramos verbos
que solo pueden tener empleo auxiliar, como \forma{soler}; luego, otros como
\forma{ir} o \forma{deber}, que funcionan como auxiliares, pero también retienen
su uso con valor pleno, y, en el polo opuesto, se encuentran ciertos verbos
plenos que comienzan a mostrar indicios de desemantización. Así, por ejemplo,
\forma{prometer} se comporta como verbo pleno en \forma{El deportista promete
llegar temprano a los entrenamientos}, donde la unión del verbo conjugado y la
forma no finita no constituye perífrasis porque admite conmutación estructural:
\forma{El deportista promete que llegará temprano}. Sin embargo, en \forma{Este
deportista promete ser la estrella del equipo}, el mismo verbo se muestra en vías
de auxiliarización, ya que no expresa ninguna promesa por parte del deportista,
sino que ha adquirido un nuevo significado por el cual el hablante expresa una
deducción anticipatoria sobre el lugar de prominencia que posiblemente alcance el
deportista en el futuro. En este segundo caso, la oración se interpreta como
\sig{Este deportista posiblemente será la estrella del equipo}.

Considerado en perspectiva histórica, el proceso por el que un verbo pleno se
convierte en un auxiliar de PV puede analizarse de la siguiente manera. Tomemos,
por ejemplo, un verbo originariamente léxico como \forma{volver}, que supone un
\sig{desplazamiento físico de retorno a un lugar}, lo que se manifiesta por el
verbo seguido de un SP con valor de dirección, como en \forma{Volvió a la
escuela}. Paulatinamente, al comenzar a emplearse con un infinitivo como
complemento de la preposición, ha ido sumando una nueva interpretación, como en
\forma{María vuelve a leer el libro}, donde es posible entender tanto que
\sig{María vuelve / regresa a un lugar determinado para leer el libro}
(\forma{María volvió a la biblioteca para leer el libro}) como que la reiteración
no se aplica al movimiento de regreso, sino a la repetición del evento
(\forma{María leyó el libro nuevamente}). Cuando se produce una nueva lectura, la
consolidación de esa interpretación tiene como contraparte la rigidización de la
estructura sintáctica. Hay una nueva correspondencia entre la forma y el
significado, y la estructura queda cristalizada como \textit{verbo auxiliar +
verbo auxiliado}. En otras palabras, mientras el verbo se comporte como el núcleo
léxico-sintáctico de la predicación, posee mayor amplitud combinatoria:
\forma{María vuelve \{a la escuela / a leer / del extranjero\}}; en cambio,
cuando se convierte en auxiliar, la estructura se vuelve más rígida y solo acepta
ir seguida del infinitivo: \forma{El árbol volvió a \{florecer / perder las
hojas\}}. Por tanto, si bien hay ciertas perífrasis consolidadas, en otros
casos, según hemos visto, puede producirse ambigüedad y, en otros, puede tratarse
de un cambio aún en proceso. De ahí que resulte clave el análisis que realicemos
al enfrentarnos a una secuencia de verbo finito + verbo no finito, ya que no solo
las distintas combinaciones, según las pruebas que acepten, pueden ser o no
perífrasis, sino que un mismo verbo, según la interpretación que habilite el
contexto, puede tener una u otra lectura. Y aún más, como vimos con
\forma{prometer}, puede haber verbos que en su sentido más general no admitan las
pruebas de auxiliarización, pero que en algunos usos «más innovadores» comiencen
a mostrar cierto desplazamiento de su sentido original.
\end{recuadro}

\begin{ejercicio}\label{ej:11-corpes}
Considere las combinaciones verbales subrayadas en las siguientes oraciones y
resuelva las consignas según se muestra en el ejemplo resuelto~1.

\begin{consignas}
  \item Determine si las secuencias que conforman constituyen o no una perífrasis
        verbal.
  \item En cada caso, fundamente su respuesta.
\end{consignas}

\begin{exe}
  \ex Alguien le comentó alguna vez que la salida más fácil para terminar una
      historia \subrayado{era matando} al protagonista.
  \ex Ves que tu hermana no \subrayado{era descabellada} del todo.
  \ex Me pareció que el Ojo \subrayado{era sacudido} por un escalofrío.
  \ex El Bando Azul de Matanzas aparece como fundado en 1910 y vinculado al
      cabildo arará. Según se infiere de los datos compilados, \subrayado{comenzó
      siendo} un coro o bando de clave y derivó a uno de guaguancó.
  \ex Él es el verdadero autor de su propia gloria y uno no puede atribuirse más
      que la alegría, pero sería egoísta si no reconociera que yo lo ayudé a
      nacer. Cuando \subrayado{comenzó pude adivinar} que tenía enormes
      condiciones. Me maravilló su actitud iconoclasta.
  \ex Lo destapó, \subrayado{comenzó a beberlo} directamente de la botella y
      apagó el monitor de vigilancia.
\end{exe}

\noindent{\footnotesize Textos extraídos de CORPES XXI.}
\pendiente{Los subrayados no se conservan en el original; se reconstruyeron a
partir del solucionario}

\begin{resuelto}
1. \forma{Era matando}: no se trata de una PV porque el gerundio puede ser
reemplazado por el adverbio \forma{así}.
\end{resuelto}
\end{ejercicio}

\section{Contenido gramatical de las perífrasis verbales}\label{sec:11-contenido}

Según adelantamos, más allá de la conformación estructural que puedan adoptar,
las perífrasis verbales también se clasifican teniendo en cuenta el aporte de
significado del auxiliar, que manifiesta valores específicos vinculados ya sea
con el tiempo \exref{ex:11-valores-a}, el modo \exref{ex:11-valores-b}, el
aspecto \exref{ex:11-valores-c} o bien con la diátesis, que es la categoría que
se relaciona con la configuración que adoptan los argumentos del verbo en la
oración \exref{ex:11-valores-d}.

\begin{exe}
  \ex
    \begin{xlist}
      \ex\label{ex:11-valores-a} María \textit{va a leer} un libro en el verano.
          $\rightarrow$ PV de temporalidad: indica futuridad.
      \ex\label{ex:11-valores-b} María \textit{puede leer} un libro durante el
          verano. $\rightarrow$ PV de modalidad: indica posibilidad.
      \ex\label{ex:11-valores-c} María \textit{comenzó a leer} un libro durante
          el verano. $\rightarrow$ PV de aspectualidad: indica inicio.
      \ex\label{ex:11-valores-d} Un libro \textit{fue leído} por María durante el
          verano. $\rightarrow$ PV de diátesis: indica voz pasiva.
    \end{xlist}
\end{exe}

En función de los distintos valores que pueden manifestar, presentamos cuatro
grupos principales de perífrasis verbales: de diátesis, aspectuales, temporales y
modales.\footnote{La clasificación de las PPVV que presentamos, basada en
\textcite{giammatteo2009marcovecchio,giammatteo2010marcovecchio}, no pretende
ser exhaustiva, sino dar cuenta de las PPVV más representativas de los distintos
valores semánticos manifestados por estas unidades.}

\subsection{Perífrasis verbal de diátesis: la voz pasiva}\label{sec:11-pasiva}

Como se vio en el \capref{7}, los verbos determinan la estructura argumental del
predicado, es decir, seleccionan los argumentos y les asignan los papeles
temáticos que les corresponden. Sin embargo, la forma en que se manifiestan estos
argumentos puede variar. En gramática recibe el nombre de \termino{diátesis} la
categoría que

\begin{cita}
establece (i) qué argumentos verbales se seleccionan para una predicación, y
(ii) qué funciones sintácticas le corresponden a cada argumento. Así,
dependiendo de las posibles formas de disponer los argumentos del verbo en una
oración, se considera que existen \textit{distintas formas de diátesis}; puede
decirse que cada diátesis constituye un patrón de correspondencia específico
entre los argumentos de un verbo y un conjunto de funciones sintácticas
\parencite[222]{munoz2018}.
\end{cita}

Para manifestar las oposiciones de diátesis, las lenguas acuden a recursos
morfológicos (flexión y derivación) o sintácticos (fundamentalmente,
perífrasis). Estas diferencias formales que expresa la diátesis, como, por
ejemplo, \forma{despertó} frente a \forma{fue despertado}, se denominan
\termino{voces} \parencite{rae2025}. Así, tanto \exref{ex:11-vidrios-a} como
\exref{ex:11-vidrios-b} expresan los mismos argumentos, pero mediante
configuraciones oracionales diferentes:

\begin{exe}
  \ex
    \begin{xlist}
      \ex\label{ex:11-vidrios-a} Juan rompió los vidrios. $\rightarrow$ Voz
          activa.
      \ex\label{ex:11-vidrios-b} Los vidrios fueron rotos (por Juan).
          $\rightarrow$ Voz pasiva.
    \end{xlist}
\end{exe}

El cambio que manifiestan \exref{ex:11-vidrios-a} y \exref{ex:11-vidrios-b} se
conoce como \termino{alternancia activa-pasiva}. Mientras la voz activa se
expresa por la forma simple del verbo \exref{ex:11-vidrios-a}, la pasiva lo hace
mediante la perífrasis formada por el auxiliar \forma{ser} y el participio del
verbo principal \exref{ex:11-vidrios-b}. El participio de la voz pasiva
concuerda en género y número con el sujeto.

\begin{exe}
  \ex
    \begin{xlist}
      \ex Un investigador de la universidad \textit{fue premiado} por su
          trayectoria académica.
      \ex Varias investigadoras de la universidad \textit{fueron premiadas} por
          sus trayectorias académicas.
    \end{xlist}
\end{exe}

La voz activa es la forma no marcada característica de los verbos transitivos.
Solo los verbos transitivos, que poseen dos argumentos ---un agente, que funciona
como sujeto, y un tema o paciente, en la función de COD---, permiten la
alternancia activa-pasiva. La forma pasiva reestructura la configuración
oracional promoviendo el tema o paciente a la función oracional más prominente,
la de sujeto, al mismo tiempo que desplaza el agente a la función marginal y
omisible de complemento agente.

Además de la pasiva perifrástica, en español también tenemos una pasiva
cuasirrefleja, formada por el clítico \forma{se} seguido de la forma activa del
verbo. La pasiva cuasirrefleja también promueve el tema o paciente a la función
de sujeto, pero no admite la manifestación sintáctica del agente.

\begin{exe}
  \ex \textit{Se rompieron} los vidrios (*por Juan).\footnote{En el \capref{10},
      §\ref{sec:10-pasivas}, se presentaron comparativamente las diferencias
      entre la pasiva perifrástica y la cuasirrefleja pasiva.}\pendiente{El
      original remite al «capítulo 9»}
\end{exe}

Asimismo, también existe otra perífrasis de diátesis, que «expresa un estado
resultante de una acción pasiva o media\ldots» \parencite[3429]{yllera1999}. Se
forma con el verbo \forma{estar} + participio concordado con el sujeto y puede
incluir un complemento agente\footnote{La adscripción de \forma{estar} como
auxiliar de perífrasis no es unívoca, dado que en este verbo conviven el
comportamiento de auxiliar, como se muestra en los ejemplos
\exref{ex:11-escombros} a \exref{ex:11-plantel}, y el de verbo copulativo, como
en \forma{Está angustiado por todo lo que le pasó}, donde el participio puede
ser conmutado por un adjetivo: \forma{Está triste por todo lo que le
pasó}.} referido al participio, que tiene interpretación pasiva. Según sugieren
\textcite[105]{garciafernandez2019}, a diferencia de las pasivas con
\forma{ser}, estas perífrasis de resultado que se construyen con \forma{estar}
«se pueden formar regularmente sobre verbos intransitivos, especialmente sobre
aquellos que tienen un incremento clítico e interpretación
media».\pendiente{El original cita «Fernández García y Krivochen»}

\begin{exe}
  \ex\label{ex:11-escombros} Los escombros \textit{estaban esparcidos} por toda
      la zona.
  \ex La reunión \textit{estuvo presidida} por el actual jefe de gobierno.
  \ex\label{ex:11-plantel} Ud. ya \textit{está incorporado} en el nuevo plantel
      docente.
\end{exe}

\subsection{Perífrasis verbales aspectuales}\label{sec:11-aspectuales}

Las perífrasis aspectuales tienen que ver con «la configuración temporal interna
(o contorno) de las situaciones (acciones, sucesos, estados, etc.)»
\parencite[346]{lyons1997}. Se reconocen dos subgrupos: las que se refieren al
desarrollo del evento y dan cuenta de las fases de su desenvolvimiento, y las que
se vinculan con su cuantificación y consideran la pluralización de los sucesos o
de las ocasiones en que ocurren, o bien su intensificación. Veamos más
detenidamente cada subtipo.

\paragraph{Perífrasis aspectuales de fase.} Manifiestan las etapas o «fases» de
desarrollo interno del evento e incluyen:

\begin{itemize}
  \item Las \termino{perífrasis incoativas}, que señalan el inicio del evento
        ---\forma{empezar / comenzar / ponerse / echar / romper a} +
        infinitivo--- o su inminencia ---\forma{estar a punto de / estar por} +
        infinitivo---.
\end{itemize}

\begin{exe}
  \ex
    \begin{xlist}
      \ex\label{ex:11-inc-a} María \textit{empezó / comenzó / se puso a
          estudiar}.
      \ex\label{ex:11-inc-b} María \textit{echó / rompió a llorar / a correr /
          *a estudiar / *a pensar}.
    \end{xlist}
  \ex\label{ex:11-parir} María \textit{estaba a punto de / por parir} y llamaba
      la atención su serenidad.
\end{exe}

Mientras que las perífrasis de \exref{ex:11-inc-a} solo señalan el comienzo del
evento de \forma{estudiar}, las de \exref{ex:11-inc-b} añaden que este se produce
de manera abrupta y repentina; por lo tanto, estas perífrasis tienen más
restricciones para combinarse con distintos verbos. Las perífrasis verbales de
\exref{ex:11-parir} describen una fase previa al inicio efectivo del evento.

\begin{itemize}
  \item Las \termino{perífrasis terminativas} indican la conclusión efectiva del
        evento ---\forma{terminar / concluir / acabar de} + infinitivo--- o su
        interrupción ---\forma{cesar / dejar de} + infinitivo---.
\end{itemize}

\begin{exe}
  \ex
    \begin{xlist}
      \ex\label{ex:11-term-a} María \textit{terminó / concluyó / acabó de
          estudiar} Derecho.
      \ex\label{ex:11-term-b} María \textit{dejó de estudiar} Derecho, pero
          retomará el año próximo.
    \end{xlist}
\end{exe}

Mientras las perífrasis de \exref{ex:11-term-a} indican la terminación del
evento, en \exref{ex:11-term-b} se señala que se suspendió la realización, sin
que ello implique que el hecho haya concluido.

\begin{itemize}
  \item Las \termino{perífrasis continuativas} marcan la continuación de un
        evento iniciado con anterioridad ---\forma{seguir / continuar} +
        gerundio---.
\end{itemize}

\begin{exe}
  \ex María \textit{siguió / continuó estudiando}.
\end{exe}

\paragraph{Perífrasis aspectuales de cuantificación.} Manifiestan un valor
incremental referido al desarrollo del evento o a su repetición, y forman dos
subgrupos:

\begin{itemize}
  \item Las \termino{perífrasis progresivas} marcan el progreso efectivo
        ---\forma{estar / andar} + gerundio--- o gradual ---\forma{ir} +
        gerundio--- del desarrollo de un evento ya comenzado.
\end{itemize}

\begin{exe}
  \ex
    \begin{xlist}
      \ex\label{ex:11-prog-a} María \textit{estaba recorriendo} el territorio. (=
          \sig{Recorría efectivamente el territorio})
      \ex\label{ex:11-prog-b} María \textit{iba recorriendo} el territorio. (=
          \sig{Recorría gradualmente el territorio})
    \end{xlist}
\end{exe}

Mientras en \exref{ex:11-prog-a} se describe un evento en su desarrollo
continuado, en \exref{ex:11-prog-b} se indica que el evento se realiza de manera
gradual.

\begin{itemize}
  \item Las \termino{perífrasis reiterativas} indican la repetición
        ---\forma{volver a} + infinitivo--- o el desarrollo interrumpido de uno o
        varios eventos ---\forma{andar} + gerundio---.
\end{itemize}

\begin{exe}
  \ex
    \begin{xlist}
      \ex\label{ex:11-reit-a} María \textit{volvió a recorrer} el territorio. (=
          \sig{Recorrió nuevamente el territorio})
      \ex\label{ex:11-reit-b} María \textit{anda recorriendo} el territorio. (=
          \sig{Realiza una recorrida que se extiende durante un cierto tiempo o
          recorre frecuentemente el territorio})
    \end{xlist}
\end{exe}

Mientras que \exref{ex:11-reit-a} indica la repetición de un evento que ya ha
concluido (\forma{volvió a leer el libro} $\rightarrow$ \sig{lo había leído y lo
leyó nuevamente}), \exref{ex:11-reit-b} «describe situaciones que se desarrollan
con interrupciones o de modo intermitente» \parencite[§\,28.14a]{rae2025},
aunque pueden involucrar un solo proceso ---\forma{anda terminando la tesis}--- o
varios ---\forma{andan buscando al sospechoso por todos lados}---.\pendiente{El
original remite a (47a) y (47b)}

\subsection{Perífrasis verbales temporales}\label{sec:11-temporales}

Dentro de las perífrasis verbales temporales incluimos la que forma los tiempos
compuestos y la que indica futuridad.

\paragraph{Perífrasis de tiempos compuestos.} Se forman con \forma{haber} +
participio no concordado. Se consideran témporo-aspectuales porque señalan un
valor temporal de anterioridad respecto de otro evento y un valor aspectual
perfectivo, que indica \sig{completamiento del evento}. Consideremos los
siguientes ejemplos:

\begin{exe}
  \ex
    \begin{xlist}
      \ex\label{ex:11-llegue-a} Cuando yo \textit{llegué}, Juan \textit{se
          retiró}.
      \ex\label{ex:11-llegue-b} Cuando yo \textit{llegué}, Juan ya \textit{se
          había retirado}.
    \end{xlist}
\end{exe}

Mientras que en \exref{ex:11-llegue-a} las formas verbales están marcando un
valor temporal de simultaneidad entre ambos eventos ---\forma{llegué} y \forma{se
retiró}---, en \exref{ex:11-llegue-b}, en cambio, lo que se está señalando es que
el retiro de Juan fue anterior a mi llegada. Además de marcar anterioridad con
respecto a otro tiempo, las perífrasis de tiempos compuestos señalan
\termino{aspecto perfecto}. \textit{Perfecto} viene del latín
\textit{perfectum}, que significa \sig{acabado o terminado}, lo que quiere decir
que, aspectualmente, los tiempos compuestos refieren un evento ya concluido y del
cual se muestran los resultados \parencite{garciafernandez1998}.

En las perífrasis que conforman los tiempos compuestos, la morfología del verbo
\forma{haber} reproduce el paradigma flexivo de los tiempos simples y constituye
formas verbales que marcan anterioridad a otro tiempo de referencia.

\begin{exe}
  \ex
    \begin{xlist}
      \ex\label{ex:11-comp-a} Ya les \textit{han dicho} la verdad. (anterioridad
          al presente $\rightarrow$ \sig{antes de ahora})
      \ex\label{ex:11-comp-b} Cuando \textit{sonó} la alarma, los ladrones ya se
          \textit{habían escapado}. (anterioridad a un evento del pasado)
      \ex\label{ex:11-comp-c} Cuando María \textit{llegue},\footnote{Después de
          \forma{cuando}, el presente de subjuntivo \forma{llegue} tiene valor
          futuro.} la fiesta ya \textit{habrá terminado}. (anterioridad a un
          evento futuro)
    \end{xlist}
\end{exe}

El pretérito perfecto compuesto \forma{han dicho} señala un evento anterior al
presente \exref{ex:11-comp-a};\footnote{En el español de la Argentina, el
pretérito perfecto también admite uso deíctico equivalente al perfecto simple:
\forma{Ayer he ido al cine} (= \forma{Ayer fui al cine}).} el pretérito
pluscuamperfecto \forma{se habían escapado} \exref{ex:11-comp-b} señala un
tiempo anterior a otro del pasado, \forma{cuando sonó la
alarma},\footnote{En el español actual, el pretérito anterior (\forma{hubo
entrado}) está prácticamente desaparecido del habla coloquial y muy restringido
en la lengua escrita \parencite{rojo1999}.} y el futuro perfecto \forma{habrá
terminado} \exref{ex:11-comp-c} indica anterioridad a un hecho futuro, \forma{la
llegada de María}.

\paragraph{Perífrasis de futuro.} Se forman con \forma{ir a} + infinitivo.
Originariamente indicaban posterioridad inmediata, es decir, un futuro próximo,
pero el sentido de inmediatez se fue perdiendo y la perífrasis pasó a significar
futuridad en general, tanto próxima \exref{ex:11-fut-a} como remota
\exref{ex:11-fut-b}.

\begin{exe}
  \ex
    \begin{xlist}
      \ex\label{ex:11-fut-a} En cinco minutos \textit{vamos a salir} para el
          campo.
      \ex\label{ex:11-fut-b} El año próximo \textit{vamos a veranear} en Brasil.
    \end{xlist}
\end{exe}

\subsection{Perífrasis verbales de modalidad}\label{sec:11-modales}

La \termino{modalidad} es la categoría que señala la actitud del hablante con
respecto a su enunciado.\footnote{Este tema se desarrolla en el
\capref{14}.\pendiente{El original remite al «capítulo 12»}}
\textcite{palmer1986} distingue entre \termino{modalidad deóntica} y
\termino{modalidad epistémica} como dos valores presentes en todas las lenguas.
El primero está vinculado con distintos matices dentro del ámbito de la
obligación y el segundo, con valores ligados a la posibilidad. Ambas modalidades
pueden ser expresadas por perífrasis verbales formadas por los auxiliares
\forma{poder, deber, tener que} + infinitivo, mientras que \forma{haber que} +
infinitivo solo admite valor deóntico.

\begin{itemize}
  \item Las \termino{perífrasis deónticas} están orientadas hacia el sujeto de la
        oración. En ellas, \forma{poder} + infinitivo adopta valor de
        \sig{capacidad} \exref{ex:11-correr} o de \sig{permiso}
        \exref{ex:11-volver}; y \forma{deber} o \forma{tener / haber que} +
        infinitivo manifiestan \sig{obligación} \exref{ex:11-bertha} o
        \sig{necesidad} \exref{ex:11-agua}.
\end{itemize}

\begin{exe}
  \ex\label{ex:11-correr} Juan \textit{puede correr} los 100 metros en 9
      segundos.
  \ex\label{ex:11-volver} María \textit{puede volver} a casa a la una.
  \ex\label{ex:11-bertha} Bertha \textit{debe entregar} la solicitud antes de las
      11.
  \ex\label{ex:11-agua} \textit{Tenemos que beber} más agua para estar sanos.
\end{exe}

\begin{itemize}
  \item Las \termino{perífrasis epistémicas} utilizan los mismos auxiliares que
        las deónticas (excepto \forma{haber que}), pero con significados
        relacionados con la \sig{posibilidad} o la \sig{probabilidad} de la
        ocurrencia del evento.
\end{itemize}

\begin{exe}
  \ex María \textit{puede llegar} a las 8~h. $\rightarrow$ Es posible que María
      llegue a las 8~h.
  \ex María \textit{debe (de) llegar} cerca de las 8~h. $\rightarrow$ Supongo que
      María llegará cerca de las 8~h.
\end{exe}

Las perífrasis modales presentan algunas particularidades. Según dijimos en
§\ref{sec:11-auxiliares}, dado que en las PPVV auxiliar y auxiliado forman una
unidad compleja, el auxiliar no admite ser negado en forma independiente del
verbo principal:

\begin{exe}
  \ex[*]{Mariana no va a no ir a la facultad.}
\end{exe}

Sin embargo, las perífrasis de modalidad (salvo \forma{tener que}), además de la
negación del auxiliar \exref{ex:11-negmod-a}, también permiten la de ambos
verbos \exref{ex:11-negmod-b}, con lectura predominantemente deóntica; en
cambio, cuando se niega solo el auxiliado, la interpretación resulta epistémica
\exref{ex:11-negmod-c}. Estas posibilidades muestran que la unión entre auxiliar
y auxiliado es más laxa en estas perífrasis verbales.

\begin{exe}
  \ex
    \begin{xlist}
      \ex\label{ex:11-negmod-a} \{\textit{No puede / no debe}\} \textit{venir}
          temprano. (No está autorizado).
      \ex\label{ex:11-negmod-b} \{\textit{No puede / no debe}\} no \textit{venir}
          temprano. (Está obligado a venir temprano).
      \ex\label{ex:11-negmod-c} \{\textit{Puede / debe}\} no \textit{venir}
          temprano. (Tiene la posibilidad de no venir / Supongo que no vendrá
          temprano).
    \end{xlist}
\end{exe}

Como vimos en \exref{ex:11-negmod-c}, las perífrasis epistémicas admiten la
negación del auxiliado \exref{ex:11-camino1}, pero cuando la negación precede a
\forma{poder}, este se interpreta como \sig{es imposible}
\exref{ex:11-camino2}:

\begin{exe}
  \ex\label{ex:11-camino1} Pudo no encontrar el camino. $\rightarrow$ Es posible
      que no haya encontrado el camino / Sospecho que no ha encontrado el camino.
  \ex\label{ex:11-camino2} No pudo perderse en ese camino. $\rightarrow$ Es
      imposible / No creo que se haya perdido en ese camino.
\end{exe}

\begin{recuadro}[Para ampliar: las perífrasis de habitualidad en la intersección entre modalidad y aspectualidad]
La perífrasis verbal de habitualidad se forma con \forma{soler, acostumbrar} y,
en el Río de la Plata, también con \forma{saber} + infinitivo. En el caso de
\forma{soler}, presenta la particularidad de que nunca actúa como verbo pleno,
sino que únicamente funciona como auxiliar seguido de infinitivo.

\begin{exe}
  \ex [L]as águilas harpías son aves muy escurridizas, que \textit{suelen volar}
      bajo entre las copas de los árboles y son difíciles de avistar. (\textit{El
      Ciudadano}, 17/02/2025)
  \ex Es por esto que te digo que no \textit{acostumbro a plantar} dalias en los
      macizos del jardín. (\textit{Revista Mi Jardín}, 2020)
  \ex Ya de vuelta de esas tierras, cuando me encontraba en la proximidad de los
      puertos, me \textit{sabía venir} la tentación de interrogar a los marinos
      que volvían de viaje. (J.~J.~Saer, \textit{El entenado})
\end{exe}

Si bien las perífrasis de habitualidad generalmente se han ubicado entre las
aspectuales, según la \textit{NGLE} \parencite[§\,28.9b]{rae2025}, «\textit{soler}
+ infinitivo ocupa un lugar intermedio entre las perífrasis modales y las
aspectuales». Anteriormente, también \textcite[40]{comrie1985} había sostenido
que «la expresión gramatical de la habitualidad siempre estará integrada en el
sistema aspectual o modal de una lengua». Pero ¿a qué se debe esta
particularidad? El valor aspectual de estas perífrasis se relaciona con que hacen
referencia a la reiteración de un evento a lo largo de un período en el que el
suceso típicamente ocurre \parencite{olbertz1998}. En cuanto a su valor modal,
deriva de que dicho evento debe cumplirse en la mayoría, pero no necesariamente
en todas las ocasiones, lo que implica que su ocurrencia tiene una cierta
regularidad en el período involucrado, aunque la reiteración puede no ser
siempre efectiva.

\begin{exe}
  \ex Juan \textit{suele llegar} (frecuentemente llega) a estas horas, pero hoy
      todavía no ha llegado (aún no está aquí).
  \ex Don Gómez \textit{sabe venir} por acá los sábados, pero este fin de semana
      no lo hemos visto.
\end{exe}

En \textcite{giammatteo2010marcovecchio} y \textcite{giammatteo2011} se plantea
que este carácter virtual de la reiteración de un evento vincula habitualidad y
posibilidad. De este modo, \forma{poder / deber (de) / tener que}, por un lado, y
\forma{soler / saber / acostumbrar}, por el otro, manifiestan grados dentro de
la modalidad epistémica. Así, \forma{soler} presupone la realización efectiva
del evento, aunque le permite al hablante evitar un compromiso absoluto con el
valor de verdad de su aseveración para todas las ocasiones; \forma{poder}, en
cambio, indica un grado aún menor de compromiso del hablante con sus dichos, ya
que solo hace referencia a una realización eventual o posible del evento, que no
confiere ninguna certeza a su realización.

\begin{exe}
  \ex Suele ir al cine los domingos. (Generalmente lo hace, aunque algún domingo
      puede no ir).
  \ex Puede ir al cine este domingo. (Existe una posibilidad de que el hecho
      suceda, aunque también podría no ocurrir).
\end{exe}
\end{recuadro}

A modo de recapitulación, en la tabla~\ref{tab:11-perifrasis} ofrecemos una
sistematización de las distintas clases de perífrasis reconocidas en el
capítulo.

{\footnotesize\setlength{\tabcolsep}{3pt}
\begin{longtable}{@{}>{\raggedright\arraybackslash}p{24mm}
    >{\raggedright\arraybackslash}p{30mm} >{\raggedright\arraybackslash}p{30mm}
    >{\raggedright\arraybackslash}p{36mm}@{}}
  \caption{Sistematización de las perífrasis verbales}\label{tab:11-perifrasis}\\
  \toprule
  Clase & Subclase & Descripción estructural & Ejemplos \\
  \midrule
  \endfirsthead
  \toprule
  Clase & Subclase & Descripción estructural & Ejemplos \\
  \midrule
  \endhead
  \multicolumn{4}{@{}p{126mm}}{\textbf{De diátesis}: modifican la relación entre
    los argumentos del verbo y las funciones sintácticas que se manifiestan en la
    oración.} \\
  \addlinespace
  Pasiva & --- & ⟨\textit{ser} + participio concordado⟩ & \textit{La casa fue
    comprada por un nuevo vecino.} \\
  Resultativa & --- & ⟨\textit{estar} + participio concordado⟩ & \textit{La casa
    está ocupada solo por muebles viejos.} \\
  \midrule
  \multicolumn{4}{@{}p{126mm}}{\textbf{Aspectuales}: aportan información sobre el
    desarrollo interno del evento.} \\
  \addlinespace
  De fase (I): etapas de desarrollo del evento & incoativas: marcan el inicio del
    evento & ⟨\textit{empezar / comenzar / echar / romper / ponerse a} + inf.⟩;
    ⟨\textit{estar a punto de / por} + inf.⟩ (inminencial) & \textit{Ana empezó a
    dibujar el plano. El niño rompió a llorar desconsolado. Los jueces están por
    dictar sentencia.} \\
  & terminativas: marcan el fin o la interrupción del evento &
    ⟨\textit{terminar / acabar / concluir de} + inf.⟩; ⟨\textit{llegar /
    alcanzar a} + inf.⟩; ⟨\textit{cesar / dejar de} + inf.⟩ (cesativa) &
    \textit{Pedro terminó de pintar la pared. Raúl dejó de leer el libro.} \\
  & continuativas: marcan la continuación de un evento en desarrollo &
    ⟨\textit{continuar / hallarse / llevar / seguir / venir} +
    ger.⟩\textsuperscript{a} & \textit{Los alumnos se hallan buscando
    información. El gato continúa persiguiendo al ratón.} \\
  De cuantificación (II): incremento en el desarrollo del evento & reiterativas:
    marcan la repetición del evento & ⟨\textit{volver a} + inf.⟩;
    ⟨\textit{andar} + ger.⟩ (reiterativa-frecuentativa) & \textit{Los libros
    volvieron a ser disfrutados por todos. Mariana anda bostezando por todas
    partes.} \\
  & progresivas: marcan el desarrollo ininterrumpido de uno o varios eventos &
    ⟨\textit{estar} + ger.⟩; ⟨\textit{ir} + ger.⟩ (progresiva gradual) &
    \textit{Margarita está traduciendo un libro al inglés. Juanita va aprendiendo
    a contar.} \\
  \midrule
  \multicolumn{4}{@{}p{126mm}}{\textbf{Temporales}: aportan información sobre el
    tiempo del evento.} \\
  \addlinespace
  De tiempos compuestos & expresan anterioridad a otro evento e indican aspecto
    perfecto & ⟨\textit{haber} + part.⟩ & \textit{Todos hemos aprobado
    Gramática.} \\
  De futuro & expresan posterioridad al tiempo presente & ⟨\textit{ir a} + inf.⟩
    & \textit{El próximo mes voy a disfrutar de mis vacaciones.} \\
  \midrule
  \multicolumn{4}{@{}p{126mm}}{\textbf{Modales}: manifiestan información sobre la
    actitud del hablante ante el contenido de su enunciado.} \\
  \addlinespace
  Deónticas & expresan obligación, necesidad o capacidad del sujeto de la oración
    & ⟨\textit{deber (de) / poder / tener que / haber que} + inf.⟩ &
    \textit{Deben estudiar este tema para el examen.} \\
  Epistémicas & expresan el grado de posibilidad que el hablante adjudica a la
    ocurrencia de un evento\textsuperscript{b} & (los mismos auxiliares, salvo
    \textit{haber que}) & \textit{Deben (de) haber estudiado mucho porque
    aprobaron.} \\
  De habitualidad (intersección entre aspecto y modalidad) & señalan que en un
    período extendido de tiempo algún suceso típicamente ocurre &
    ⟨\textit{soler} + inf.⟩; ⟨\textit{acostumbrar a} + inf.⟩; ⟨\textit{saber}
    + inf.⟩ & \textit{En el pueblo suelen festejar el 25 de Mayo con locro y
    empanadas. Sabe estar siempre en su casa tomando mate.} \\
  \bottomrule
\end{longtable}}

{\scriptsize\noindent\textsuperscript{a} Dada la índole introductoria de este
libro, omitimos detallar los diferentes matices aspectuales que manifiestan
estas PPVV. \textsuperscript{b} Los distintos grados de posibilidad o
probabilidad de la ocurrencia de un evento dependen del grado de certeza con que
el hablante lo comunica (información de primera mano, inferencia directa,
indirecta, etc.). Estos temas se amplían en el \capref{14}, dedicado a la
modalidad.\par}

\begin{ejercicio}\label{ej:11-pleno}
En las siguientes oraciones, diferencie el significado del verbo en su uso pleno
del que aporta como auxiliar de perífrasis, según se muestra en el ejemplo
resuelto~1.

\begin{exe}
  \ex
    \begin{xlist}
      \ex Juana va todos los días a la facultad a estudiar Derecho.
      \ex Juana va a estudiar Derecho el año próximo.
    \end{xlist}
  \ex
    \begin{xlist}
      \ex El iracundo rockero rompió la guitarra en el escenario.
      \ex El iracundo rockero rompió a llorar desconsoladamente.
    \end{xlist}
  \ex
    \begin{xlist}
      \ex Al final del período cretácico, ya no había dinosaurios.
      \ex Antes de su muerte habían ocurrido algunos hechos sospechosos.
    \end{xlist}
  \ex
    \begin{xlist}
      \ex Mi abuelo era pelado.
      \ex Mi abuelo era conocido por todos en el pueblo.
    \end{xlist}
  \ex
    \begin{xlist}
      \ex El jugador está en el área de penales.
      \ex El jugador está por salir al campo de juego.
    \end{xlist}
\end{exe}

\begin{resuelto}
1a. \sig{dirigirse a}. 1b. Futuridad.
\end{resuelto}
\end{ejercicio}

\begin{ejercicio}\label{ej:11-significados}
Determine si las siguientes oraciones contienen perífrasis verbales y, en caso de
que lo hagan, señale los significados que manifiestan, según se muestra en los
ejemplos resueltos.

\begin{exe}
  \ex María necesita aprender gramática para aprobar su examen.
  \ex María tiene que aprender gramática para aprobar su examen.
  \ex María está aprendiendo gramática para aprobar su examen.
  \ex María aprende cantando.
  \ex María va aprendiendo gramática.
  \ex María continúa aprendiendo gramática.
  \ex María es dominada por su pasión por la gramática.
  \ex María está sentada en primera fila.
  \ex María suele estudiar gramática.
  \ex María empezó a estudiar gramática.
  \ex María no volvió a rendir el examen de gramática porque lo aprobó.
  \ex María no vuelve a estudiar para el examen de gramática, sino para
      descansar.
\end{exe}

\begin{resuelto}[Ejemplos resueltos]
1. No es PV. \textit{Aclaración}: en \forma{necesita aprender} no es posible
aplicar la prueba de conmutación estructural, según se explicó en
§\ref{sec:11-auxiliares}, porque el sujeto de ambos verbos es el mismo,
\forma{María}; pero la aplicación de la prueba es posible si cambiamos la persona
del sujeto del primer verbo: \forma{Yo necesito que María aprenda gramática}. 2.
PV de modalidad deóntica.
\end{resuelto}
\end{ejercicio}

\begin{recuadro}[Para ampliar: equivalencias semánticas entre perífrasis verbales y adverbios]
Algunos de los significados que expresan las perífrasis verbales pueden también
manifestarse a través de adverbios, como puede observarse en los siguientes pares
de oraciones, en los que las versiones de (a) expresan respectivamente futuridad,
continuidad y posibilidad mediante adverbios, y las de (b) lo hacen a través de
perífrasis verbales.

\begin{exe}
  \ex Futuridad:
    \begin{xlist}
      \ex \textit{Mañana} voy a Buenos Aires.
      \ex \textit{Voy a ir} a Buenos Aires.
    \end{xlist}
  \ex Continuidad:
    \begin{xlist}
      \ex \textit{Todavía} Juan está en Buenos Aires.
      \ex Juan \textit{sigue estando} en Buenos Aires.
    \end{xlist}
  \ex Posibilidad:
    \begin{xlist}
      \ex \textit{Posiblemente / Tal vez} Juan esté en Buenos Aires.
      \ex Juan \textit{puede estar} en Buenos Aires.
    \end{xlist}
\end{exe}
\end{recuadro}

\begin{ejercicio}\label{ej:11-reformular}
Reformule las siguientes oraciones con una perífrasis verbal de manera que
conserven el mismo significado, según se muestra en los ejemplos resueltos.
Tenga en cuenta los segmentos subrayados.\pendiente{Los subrayados no se
conservan en el original}

\begin{exe}
  \ex Se compraron algunas provisiones.
  \ex Tal vez te guste la nueva novela.
  \ex Mañana le doy a mi compañero la contraseña.
  \ex Le ordeno que se detenga.
  \ex Aquí nieva con frecuencia en invierno.
  \ex Posiblemente, Ricardo se ha dormido.
  \ex Enseguida inicio el armado del equipo.
  \ex El escritor recién terminó su última novela.
  \ex El escritor vivió durante muchos años en San Juan, pero ahora vive en
      Santiago de Chile.
  \ex Mariela nuevamente será mamá.
  \ex El profesor les pidió a sus alumnos que mantengan el ritmo de estudio.
  \ex Por favor, ya no peleen más.
  \ex Todos sabemos que Miguel tuvo la capacidad de vencer a su rival.
  \ex Finalmente, se dieron las condiciones para que Rosa duerma.
\end{exe}

\begin{resuelto}[Ejemplos resueltos]
1. Algunas provisiones fueron compradas. 2. La nueva novela te puede gustar.
\end{resuelto}
\end{ejercicio}

\section{Agrupaciones o cadenas perifrásticas}\label{sec:11-cadenas}

Las perífrasis verbales también pueden combinarse entre sí para formar unidades
verbales complejas que reúnen más de un auxiliar. En estos casos, solo el primer
auxiliar es una forma conjugada; los demás mantienen la forma no finita que
corresponde a la clase de perífrasis que integran ---\forma{pueden empezar a
comer, había sido designado, iban a estar conversando}---.

\begin{ejercicio}\label{ej:11-cadenas}
Observe los ejemplos siguientes. Luego realice una descripción estructural de
las perífrasis que integran cada núcleo verbal e indique el significado que
transmite cada uno de sus constituyentes, según se muestra en el ejemplo
resuelto~1.

\begin{exe}
  \ex María \textit{debe seguir estudiando} en la universidad.
  \ex María \textit{pudo terminar de estudiar} Derecho el año pasado.
  \ex María \textit{va a poder empezar a estudiar} en la universidad el próximo
      año.
  \ex Juan y María \textit{deberían haber empezado a estudiar} para los exámenes
      hace rato.
  \ex Juan \textit{tiene que poder dejar de comer} grasas saturadas.
\end{exe}

\begin{resuelto}
1. \forma{Debe} (obligatoriedad) + \forma{seguir} (continuidad) + gerundio
\forma{estudiando}.
\end{resuelto}
\end{ejercicio}

\subsection{Ámbito oracional de las perífrasis verbales}\label{sec:11-ambito}

Ahora bien, ¿qué sucede con el orden entre perífrasis verbales cuando se
combinan entre sí formando cadenas? Observemos los siguientes casos:

\begin{exe}
  \ex\label{ex:11-chomsky}
    \begin{xlist}
      \ex[*]{\label{ex:11-chomsky-a} El libro de Chomsky \textit{empezó a haber
          sido leído} por los estudiantes.}
      \ex\label{ex:11-chomsky-b} El libro de Chomsky \textit{pudo ser leído} por
          los estudiantes.
      \ex[*]{\label{ex:11-chomsky-c} El libro \textit{fue debido de ser leído}
          por los estudiantes.}
      \ex\label{ex:11-chomsky-d} Los estudiantes \textit{habían vuelto a leer} el
          libro de Chomsky antes del examen.
      \ex[*]{\label{ex:11-chomsky-e} Los estudiantes \textit{continuaron habiendo
          comprado} libros.}
    \end{xlist}
\end{exe}

Según vemos en los ejemplos de arriba, algunos órdenes entre PPVV son posibles,
mientras que otros producen resultados agramaticales. Pero ¿qué determina la
gramaticalidad o agramaticalidad de las oraciones anteriores? Aunque podría
suponerse que la combinación de perífrasis verbales se da en un orden aleatorio,
su ordenamiento secuencial tiene que ver con los niveles o capas de la estructura
oracional que hemos visto en el \capref{6}.

\begin{destacado}
\textbf{Para recordar.} Según la propuesta de \textcite{rizzi1995}, la oración se
organiza en tres capas o niveles (\textit{shells}): la \termino{capa léxica} o
del SV, donde se produce la vinculación del V con sus argumentos y
modificadores; la \termino{capa oracional} o del SF, que pone en relación el
sujeto con el predicado y en la que la oración se temporaliza, y la
\termino{capa informativa} o del SComp, que incluye la perspectiva del hablante y
establece la articulación con el nivel del discurso.
\end{destacado}

Dentro del dominio del SV se ubica la perífrasis de voz pasiva,\pendiente{El
original dice «voz activa»} que es la más interna y afecta a la estructura
argumental del verbo.

\begin{exe}
  \ex {}[\etq{PAS} Fue reconocido] por su obra literaria.
\end{exe}

También dentro del mismo nivel del SV se sitúan, en orden correlativo, las
perífrasis aspectuales de fase, que preceden a la pasiva \exref{ex:11-asp-a};
luego, las cuantificacionales, que anteceden a las de fase y a la pasiva
\exref{ex:11-asp-b}, y, por último, las de modalidad deóntica, que preceden a
todas las anteriores, \exref{ex:11-deont-a} y \exref{ex:11-deont-b}:

\begin{exe}
  \ex
    \begin{xlist}
      \ex\label{ex:11-asp-a} {}[\etq{ASP I} Empezó a [\etq{PAS} ser reconocido]]
          por su obra literaria en años recientes.
      \ex\label{ex:11-asp-b} {}[\etq{ASP II} Está [\etq{ASP I} empezando a
          [\etq{PAS} ser reconocido]]] por su obra literaria recién ahora.
    \end{xlist}
  \ex
    \begin{xlist}
      \ex\label{ex:11-deont-a} {}[\etq{DEÓNT} Debe [\etq{ASP I} empezar a
          [\etq{PAS} ser premiado]]] por su obra literaria.
      \ex\label{ex:11-deont-b} {}[\etq{DEÓNT} Puede [\etq{ASP II} volver a
          inscribirse]] en el curso, si completa este formulario.
          (\forma{Puede} = \sig{se le permite}).
    \end{xlist}
\end{exe}

En el nivel del SF se ubican las perífrasis temporales, que se vinculan con el
tiempo gramatical de la oración. Estas perífrasis anteceden a todas las que se
sitúan en el nivel previo del SV (pasivas, aspectuales de fase y de
cuantificación, y modales deónticas).

\begin{exe}
  \ex {}[\etq{TEMP} Había [\etq{ASP I} terminado de estudiar]] ese año.
  \ex {}[\etq{TEMP} Va a [\etq{ASP II} volver a [\etq{PAS} ser premiado]]] por su
      trayectoria.
\end{exe}

Finalmente, en la capa más externa de la oración, correspondiente al SComp,
donde se asientan los elementos que conectan el contenido proposicional de la
oración con la fuerza ilocutiva del hablante, se ubican las perífrasis de
modalidad epistémica:

\begin{exe}
  \ex {}[\etq{EPIST} Deben (de) [\etq{TEMP} haber [\etq{PAS} sido capturados]]]
      por la policía. (\forma{Deben (de)} = \sig{es posible que hayan sido
      capturados}).
  \ex Los manifestantes [\etq{EPIST} pueden [\etq{ASP II} estar [\etq{PAS} siendo
      dispersados]]]. (\forma{Pueden} = \sig{es posible que estén siendo
      dispersados}).
\end{exe}

Por tanto, según acabamos de ver, en la combinación de perífrasis el orden de los
auxiliares que se observa es el siguiente:

\begin{center}
[EPIST [TEMP [DEÓNT [ASP II (cuant.) [ASP I (de fase) [PAS]]]]]]
\end{center}

En el caso de los auxiliares de habitualidad, como \forma{soler, saber} y
\forma{acostumbrar}, una de las pruebas que se aportan en
\textcite{giammatteo2009marcovecchio,giammatteo2010marcovecchio} para demostrar
que también corresponden al ámbito del SComp es que estos auxiliares no aceptan
combinarse secuencialmente con los modales epistémicos, por lo que estarían
ocupando la misma posición estructural, lo que hace incompatible su
combinación.

\begin{exe}
  \ex
    \begin{xlist}
      \ex[*]{\textit{Suele deber de llegar} temprano.}
      \ex[*]{\textit{Puede soler venir} los viernes.}
    \end{xlist}
\end{exe}

Al igual que \forma{poder} y los demás auxiliares epistémicos, \forma{soler}
ocupa siempre la posición más externa en la estructura oracional, por lo que
antecede al resto de los auxiliares que pueden conformar perífrasis:

\begin{exe}
  \ex
    \begin{xlist}
      \ex María \textit{suele tener que trabajar} los sábados.
      \ex[*]{María \textit{tiene que soler trabajar} los sábados.}
    \end{xlist}
\end{exe}

\begin{recuadro}[Para ampliar: ámbito y posibilidades de desplazamiento]
Si bien la concatenación de perífrasis responde a restricciones no aleatorias, su
reflejo en las lenguas puede variar. Mientras en inglés, que tiene orden fijo, la
secuencia de auxiliares permanece siempre invariable, en español, lengua de orden
más libre, la secuencialidad entre auxiliares no es tan rígida y permite algunos
desplazamientos. No obstante, los movimientos permitidos se circunscriben a los
límites que imponen las capas o niveles de estructuración en los que se ubican
las distintas perífrasis. Así, aquellas que se encuentran dentro del ámbito del
SV (pasivas, aspectuales de fase, cuantificacionales y deónticas) admiten más de
una posición entre ellas.

\begin{exe}
  \ex Este artículo [\etq{PAS} fue [\etq{ASP I} empezado a / terminado de
      revisar]] la semana pasada.
  \ex Este artículo [\etq{ASP II} volvió a [\etq{DEÓNT} tener que [\etq{PAS} ser
      leído]]] entre los académicos.
  \ex {}[\etq{ASP II} Está [\etq{DEÓNT} pudiendo comprar]] por Internet
      nuevamente. (\forma{Poder} = \sig{capacidad o permiso}).
\end{exe}

Sin embargo, ni la pasiva ni las aspectuales o las deónticas pueden ocurrir fuera
del límite del SV, por lo que no pueden preceder ni a las temporales,
\exref{ex:11-desp1} y \exref{ex:11-desp2}, ni a las epistémicas, excepto que
admitan cambio de significado. Así, \forma{ir a} + infinitivo pierde su valor de
futuridad cuando se pospone a una PV aspectual o de modalidad deóntica
\exref{ex:11-desp2}, y \forma{poder, deber} y \forma{tener que}, cuando van
precedidos por perífrasis aspectuales, solo admiten interpretación deóntica
\exref{ex:11-desp3}.\pendiente{El original remite aquí a (86) y (88)}

\begin{exe}
  \ex[*]{\label{ex:11-desp1} [\etq{ASP I} Empezó a [\etq{TEMP} haber
      conseguido]] otra oportunidad.}
  \ex[*]{\label{ex:11-desp2} [\etq{ASP I/II} Volvió a / Terminó de [\etq{TEMP} ir
      a escribir]] un nuevo ensayo.}
  \ex\label{ex:11-desp3} {}[\etq{ASP I} Empezó a [\etq{DEÓNT} tener que
      resolver]] solo sus problemas.
\end{exe}

Asimismo, las perífrasis témporo-aspectuales, que se alojan en el nivel de la
capa flexiva, no tienen alcance sobre el ámbito más externo, correspondiente a la
modalidad oracional, es decir que el orden de los auxiliares permite determinar
el valor deóntico o epistémico de \forma{poder} y \forma{deber}: a) pospuestos a
la perífrasis temporal con \forma{haber}, se interpretan con valor deóntico
(dentro de la capa del SV) \exref{ex:11-ha-podido}; b) antepuestos a la
perífrasis temporal con \forma{haber}, quedan fuera del ámbito del SF y su única
interpretación posible es con valor epistémico \exref{ex:11-pudo-haber}.

\begin{exe}
  \ex\label{ex:11-ha-podido} Ella [\etq{TEMP} ha [\etq{DEÓNT} podido traducir]]
      el libro. $\rightarrow$ Lectura deóntica: \sig{Ella ha tenido la capacidad
      de traducir el libro}.
  \ex\label{ex:11-pudo-haber} Ella [\etq{EPIST} pudo [\etq{TEMP} haber traducido]]
      el libro. $\rightarrow$ Lectura epistémica: \sig{Es posible que fuera ella
      quien lo tradujo}.
\end{exe}
\end{recuadro}

\begin{ejercicio}\label{ej:11-orden}
Luego de haber reflexionado sobre el orden secuencial al que responden las
cadenas de PPVV, según se muestra en el ejemplo resuelto \exref{ex:11-chomsky-a},
¿podría decir en qué consiste la gramaticalidad o agramaticalidad de las
oraciones de \exref{ex:11-chomsky-b} a \exref{ex:11-chomsky-e}?

\begin{resuelto}
\exref{ex:11-chomsky-a}: es agramatical porque la PV de inicio con
\forma{empezar} no puede anteceder a la de tiempo compuesto.
\end{resuelto}
\end{ejercicio}

\begin{ejercicio}\label{ej:11-tabla}
Teniendo en cuenta los significados que manifiestan, ubique en la tabla de abajo
los verbos auxiliares y auxiliados de las combinaciones de perífrasis de las
oraciones siguientes, según se muestra en el ejemplo resuelto.

\begin{exe}
  \ex Los recién llegados deben comenzar a aprender las costumbres del lugar.
  \ex Woody Allen puede haber tenido que dejar de ser atendido por su
      psicoanalista.
  \ex Los alumnos suelen tener que volver a leer los textos de gramática.
  \ex Los libros han tenido que ser restaurados por especialistas.
  \ex Las páginas del libro deben (de) haber sido deterioradas por el paso del
      tiempo.
  \ex Mañana vamos a comenzar a practicar para el examen.
\end{exe}

\begin{center}
\resizebox{\linewidth}{!}{%
\begin{tabular}{@{}c *{6}{c} c@{}}
  \toprule
  & \multicolumn{6}{c}{Verbos auxiliares} & \\
  \cmidrule(lr){2-7}
  Ej. & \makecell{Mod. epistémica /\\habitualidad} & Temporalidad &
    \makecell{Mod.\\deóntica} & \makecell{Aspect.\\tipo II} &
    \makecell{Aspect.\\tipo I} & \makecell{Voz\\pasiva} & \makecell{Verbo\\auxiliado} \\
  \midrule
  1 & --- & --- & \textit{deben} & --- & \textit{comenzar a} & --- &
    \textit{aprender} \\
  2 & & & & & & & \\
  3 & & & & & & & \\
  4 & & & & & & & \\
  5 & & & & & & & \\
  6 & & & & & & & \\
  \bottomrule
\end{tabular}}
\end{center}
\end{ejercicio}

\begin{ejercicio}\label{ej:11-valores}
Identifique los valores gramaticales que manifiestan los auxiliares de las
perífrasis verbales subrayadas en las siguientes oraciones, según se muestra en
el ejemplo resuelto~1.

\begin{exe}
  \ex Los profesores \subrayado{van a tener que entregar} los programas con
      varios cambios significativos.
  \ex Tranquilamente, Susana \subrayado{puede seguir siendo considerada} la mejor
      alumna del colegio.
  \ex Jorge regresó al hospital para \subrayado{poder volver a ser atendido} por
      la enfermera.
  \ex Camilo \subrayado{ha vuelto a dejar de tomar} su medicina.
\end{exe}

\begin{resuelto}
1. Valores gramaticales: temporal de futuridad (\forma{van a}) + modal deóntico
(\forma{tener que entregar}). La PV de modalidad tiene valor deóntico porque
está precedida por la temporal.
\end{resuelto}
\end{ejercicio}

\begin{ejercicio}\label{ej:11-entrar}
Observe el funcionamiento de los complejos verbales encabezados por el verbo
\forma{entrar} en los siguientes fragmentos y luego resuelva lo pedido en las
consignas, según se muestra en el ejemplo resuelto~1.

\begin{consignas}
  \item Aplique las pruebas correspondientes para determinar si se trata de
        perífrasis verbales o no.
  \item Identifique el significado que transmite \forma{entrar} + infinitivo en
        sus usos perifrásticos.
\end{consignas}

\begin{exe}
  \ex Me levanto y voy a bañarme. \textit{Entro a trabajar} a las nueve, pero me
      he levantado media hora antes de lo acostumbrado; son las siete.
  \ex (\ldots) De allí, despertada su vocación militar por la guerra,
      \textit{entró a formar} parte del Ejército del Rin, distinguiéndose en la
      batalla de Wissembourg, donde fue herido en un brazo.
  \ex (\ldots) La astucia imperceptible de Antonio \textit{entró a funcionar}
      ---trataba de borrar el reciente traspié y congraciarse---.
  \ex Treparon una reja para \textit{entrar a robar} a una peluquería de Palermo
      y los detuvieron \textit{in fraganti}.
  \ex Luego se \textit{entró a absorber} el ahorro nacional que se volcaba al
      Mercado de Capitales. Las empresas no tenían posibilidad de competir con
      los bonos del Estado.
  \ex En un primer momento nos quedamos anonadados; unos segundos después el
      resto de los pibes se \textit{entró a reír} a carcajadas, algunas minas
      también, Adriana lució su mejor cara de perra triunfante y Alicia se apuró
      en ir a consolar a Patri.
\end{exe}

\noindent{\footnotesize Las oraciones 1--3, 5 y 6 fueron extraídas de M.~Davies
(2016--), \textit{Corpus del Español: Web/Dialects}
(\url{http://www.corpusdelespanol.org/web-dial/}), y la 4, de \textit{Infobae}
(17/06/2024).}

\begin{resuelto}
1. No es PV: admite conmutación estructural (\forma{entro y trabajo / entro para
trabajar}) y \forma{entrar} actúa como verbo pleno.
\end{resuelto}
\end{ejercicio}

\section{Locuciones verbales}\label{sec:11-locuciones}

Las \termino{locuciones verbales} son unidades complejas en las que el primer
constituyente, como sucede en las perífrasis verbales, es un verbo; pero, a
diferencia de estas últimas, el segundo constituyente puede ser tanto una forma
verbal (\forma{echar a perder}) como otra clase de palabra (\forma{tomar el
pelo, dar fe, echar de menos}). Además, en las locuciones verbales los
constituyentes pueden incluir más de dos palabras (\forma{hacer la cama} =
\sig{engañar}; \forma{dejar sin aliento} = \sig{impactar}). Otra diferencia
fundamental respecto de las perífrasis es que las locuciones verbales son
conjuntos de palabras con un significado unitario que, en líneas generales, no
deriva composicionalmente de los significados que individualmente aportan sus
constituyentes, es decir, las locuciones no pueden ser interpretadas
literalmente. Podemos ejemplificar estos conceptos con los siguientes casos:

\begin{exe}
  \ex\label{ex:11-entender} María \textit{dio a entender} sus pretensiones.
  \ex Pedro \textit{mete la pata} todo el tiempo.
  \ex La profesora \textit{echa de menos} su tierra natal.
  \ex\label{ex:11-toalla1} La estudiante \textit{tiró la toalla} con la lectura
      del libro de Chomsky.
\end{exe}

Los significados de las locuciones en \exref{ex:11-entender} a
\exref{ex:11-toalla1}, que pueden glosarse como \sig{sugerir}, \sig{equivocarse},
\sig{extrañar} y \sig{rendirse}, no se desprenden de los de las palabras que las
integran, sino, más bien, de un proceso de reinterpretación por el que llegan a
adquirir un significado unitario. El grado de integración de las unidades que
componen las locuciones se corrobora con ciertas propiedades que actúan sobre el
conjunto y no sobre cada unidad por separado. Así, como muestran los siguientes
ejemplos, los componentes no pueden ser modificados \exref{ex:11-loc-a} ni
negados \exref{ex:11-loc-b} en forma independiente, aunque admiten adverbios
adjuntos \exref{ex:11-pata-a} e interpolados \exref{ex:11-pata-b}.

\begin{exe}
  \ex
    \begin{xlist}
      \ex[*]{\label{ex:11-loc-a} María \textit{dio} hoy \textit{a conocer} ayer
          sus pretensiones.}
      \ex[*]{\label{ex:11-loc-b} María no \textit{dio a} no \textit{conocer} sus
          pretensiones.}
    \end{xlist}
  \ex
    \begin{xlist}
      \ex\label{ex:11-pata-a} Pedro \textit{mete la pata} sistemáticamente.
      \ex\label{ex:11-pata-b} Pedro \textit{mete} sistemáticamente \textit{la
          pata}.
    \end{xlist}
\end{exe}

Al igual que las perífrasis verbales, las locuciones conforman unidades
complejas que pueden funcionar como núcleo del predicado.

\begin{exe}
  \ex
    \begin{xlist}
      \ex El artista [\etq{SV/Pred} [\etq{V/n} \textit{ha terminado}] con el
          espectáculo]. $\rightarrow$ PV.
      \ex El artista [\etq{SV/Pred} [\etq{V/n} \textit{dio por terminado}] el
          espectáculo]. $\rightarrow$ LV.
    \end{xlist}
\end{exe}

\subsection{Reconocimiento}\label{sec:11-reconocimiento-lv}

A diferencia de las perífrasis verbales, que, según hemos visto, aceptan
conmutación léxica (§\ref{sec:11-lexica}), pero no estructural
(§\ref{sec:11-estructural}), las locuciones son complejos verbales más
solidificados y no admiten ninguno de estos cambios, por lo que la aplicación de
estas dos pruebas puede permitirnos distinguir entre ambas estructuras.

\begin{exe}
  \ex\label{ex:11-gobernador} Perífrasis verbales:
    \begin{xlist}
      \ex El gobernador \textit{no pudo regresar} a su provincia.
      \ex\label{ex:11-gobernador-b} Conmutación léxica: El gobernador no pudo
          volver / retornar a su provincia.
      \ex\label{ex:11-gobernador-c} Conmutación estructural: *El gobernador no
          pudo que regresaba a su provincia.
    \end{xlist}
  \ex\label{ex:11-partido} Locuciones verbales:
    \begin{xlist}
      \ex María \textit{echó a perder} el partido.
      \ex\label{ex:11-partido-b} Conmutación léxica: *María echó a ganar el
          partido.
      \ex\label{ex:11-partido-c} Conmutación estructural: *María echó y perdió el
          partido.
    \end{xlist}
\end{exe}

Según vemos, en \exref{ex:11-gobernador} es posible la conmutación léxica
\exref{ex:11-gobernador-b}, pero no la conmutación estructural
\exref{ex:11-gobernador-c}, por lo que se trata de una PV; en cambio, en
\exref{ex:11-partido}, el complejo \forma{echó a perder} no acepta ni conmutación
léxica \exref{ex:11-partido-b} ni conmutación estructural
\exref{ex:11-partido-c}, por lo que constituye una locución verbal (LV).

No obstante, hay que tener en cuenta que en la interpretación del significado de
las locuciones verbales y en su reconocimiento como tales incide el contexto en
el que se encuentran. Por ejemplo, en \exref{ex:11-toalla-a}, \forma{tiró la
toalla} admite una interpretación literal, en la que \forma{tiró} es un verbo
pleno, que actúa como núcleo del predicado y va seguido de un COD y de un CCL; en
cambio, en \exref{ex:11-toalla-b}, se trata de una locución porque ambos
constituyentes se interpretan como \sig{rendirse}. El caso de
\exref{ex:11-toalla-c} es interesante porque admitiría ambas interpretaciones: el
entrenador puede haber tirado literalmente la toalla con que le seca el sudor al
boxeador (sentido pleno de \forma{tirar}), pero en el boxeo ese gesto se
interpreta mayoritariamente como \sig{rendirse}, es decir, como una locución
verbal.

\begin{exe}
  \ex
    \begin{xlist}
      \ex\label{ex:11-toalla-a} María \textit{tiró la toalla} en el cesto de la
          ropa sucia.
      \ex\label{ex:11-toalla-b} María \textit{tiró la toalla} después de
          desaprobar el examen.
      \ex\label{ex:11-toalla-c} El entrenador del boxeador \textit{tiró la toalla}
          en mitad de la pelea.
    \end{xlist}
\end{exe}

\begin{recuadro}[Para ampliar: locuciones verbales frente a verbos de apoyo]
Los denominados \termino{verbos de apoyo} o \termino{de soporte} (\forma{dar un
paseo, hacer una pregunta, formar parte}), que son estructuras semilexicalizadas
formadas por un verbo y un SN que funciona como su complemento, poseen puntos de
contacto con las locuciones verbales. Los verbos de apoyo, que muchas veces
tienen un equivalente verbal aproximado (\forma{dar un paseo} = \forma{pasear};
\forma{hacer una pregunta} = \forma{preguntar}; \forma{formar parte} =
\forma{participar}), difieren de las locuciones en que su significado suele ser
composicional y se deduce de la suma de sus partes. Además, en el caso de los
verbos de apoyo, las estructuras que forman con el SN que los acompaña están
menos lexicalizadas, por lo que los sustantivos de las construcciones admiten
variantes: \forma{dar un paseo / una caminata / unas vueltas}. Además, a partir
del sustantivo se pueden formar oraciones subordinadas\footnote{Las oraciones
subordinadas, tanto las de relativo como las completivas, se estudian en el
\capref{13}.} de relativo con el verbo de apoyo como núcleo del predicado de la
subordinada (\forma{el paseo que di, la pregunta que hice}), y el sustantivo
admite modificadores: \forma{di un paseo muy largo; hice una pregunta bastante
difícil; formé parte fundamental del directorio de la empresa}.
\end{recuadro}

\begin{ejercicio}\label{ej:11-locuciones}
Responda las siguientes consignas, según se muestra en el ejemplo resuelto~1.

\begin{consignas}
  \item Determine si las siguientes oraciones contienen PPVV o LLVV aplicando las
        pruebas estudiadas: conmutación léxica y conmutación estructural.
  \item En el caso de las locuciones, proponga una forma verbal de significado
        equivalente.
\end{consignas}

\begin{exe}
  \ex Juan le \textit{dio una mano} a Ricardo con su nuevo negocio.
  \ex Todos \textit{hemos aprendido} mucho con los últimos acontecimientos.
  \ex María \textit{echó de menos} a sus padres durante el viaje a Mendoza.
  \ex Carlos \textit{hizo añicos} su auto deportivo.
  \ex El mes próximo \textit{vamos a disfrutar} de una nueva puesta teatral de
      \textit{Hamlet}.
  \ex El político \textit{dio a conocer} a sus votantes su programa de gobierno.
  \ex Nadie \textit{tuvo en cuenta} sus opiniones sobre esta locución.
  \ex Todos \textit{cayeron en la cuenta} de que la caída del régimen se
      aproximaba.
  \ex Los autores \textit{dan por terminado} este capítulo del libro.
  \ex Siempre \textit{puede haber} algo más para agregar.
\end{exe}

\begin{resuelto}
1. a) LV: no acepta ni conmutación léxica ni conmutación estructural. b)
Equivale a \sig{ayudar}.
\end{resuelto}
\end{ejercicio}

\begin{recuadro}[Reflexión final]
En este capítulo hemos presentado un panorama de conjunto sobre un grupo de
estructuras de las que el hablante se sirve para indicar, de modo más específico
que con las formas verbales simples, distintos matices respecto del evento: si
ocurre antes de otro suceso, si ocurrirá en el futuro, si acaba de terminar o
está próximo a suceder, si se repite, etc. Las perífrasis y locuciones verbales
estudiadas son unidades sintácticas, ya que funcionan conjuntamente como núcleo
del predicado, pero también son unidades semánticas: tienen significado
unitario. No obstante, difieren porque, en el caso de las perífrasis, el
significado se interpreta composicionalmente, como resultado de lo que expresa
cada uno de sus constituyentes, mientras que, en el caso de las locuciones, se
trata de unidades más fosilizadas cuya interpretación no es literal.
\end{recuadro}

% --- Secciones finales -------------------------------------------------
\lecturasrecomendadas

\solucionario

\begin{solucion}{\ref{ej:11-reconocer}}
3. No es PV porque indica dos eventos y admite conmutación estructural:
\forma{Mi abuela ni salió ni paseó durante la pandemia}. 4. Es ambigua y posee
dos interpretaciones: i) no es PV porque admite conmutación estructural:
\forma{Volvió a la facultad para hacer el trabajo práctico}; ii) es PV y no
admite conmutación estructural. Acepta la paráfrasis \forma{Hizo el trabajo
práctico nuevamente}. 5. Es PV porque señala un único evento y no admite
conmutación estructural: \juicio{*}{\forma{El domingo empezó que llovía desde
temprano}}.
\end{solucion}

\begin{solucion}{\ref{ej:11-series}}
a) 2. \forma{Son devoradas} es PV. Se trata de un evento único; admite
conmutación léxica, \forma{son tragadas}, pero no conmutación estructural,
\juicio{*}{\forma{son y se devoran}}. 3. \forma{Puede distraer} es PV. Se trata
de un evento único; admite conmutación léxica, \forma{puede entretener}, pero no
conmutación estructural, \juicio{*}{\forma{puede y distrae}}. 4.
\forma{Comenzarlas a ver} es PV. Se trata de un evento único; admite conmutación
léxica, \forma{comenzarlas a mirar}, pero no conmutación estructural,
\juicio{*}{\forma{comenzar y verlas}}. 5. \forma{Permiten conectarnos} no es PV.
Se trata de dos eventos; admite conmutación léxica, \forma{permiten
vincularnos}, y conmutación estructural, \forma{permiten que nos conectemos}. 6.
\forma{Suelen proveer} es PV. Se trata de un evento único; admite conmutación
léxica, \forma{suelen suministrar}, pero no conmutación estructural,
\juicio{*}{\forma{suelen y proveen}}.

b) El fragmento reescrito reemplaza las formas en cursiva por \forma{nos hacen
experimentar, son tragadas, puede entretener, comenzarlas a mirar, permiten
vincularnos} y \forma{suelen suministrar}.\pendiente{El original transcribe el
fragmento completo reescrito}
\end{solucion}

\begin{solucion}{\ref{ej:11-tato}}
2. \forma{Iremos a comprar}: es PV. DE: V auxiliar \forma{ir} + \forma{a} + V
auxiliado \forma{comprar} en infinitivo. \forma{Tener que entregar}: es PV. DE: V
auxiliar \forma{tener} + \forma{que} + V auxiliado \forma{entregar} en
infinitivo. \forma{Dicen ser}: no es PV. \forma{Están corriendo}: es PV. DE: V
auxiliar \forma{estar} + V auxiliado \forma{correr} en gerundio. 3.
\forma{Estamos viviendo}: es PV. DE: V auxiliar \forma{estar} + V auxiliado
\forma{vivir} en gerundio. \forma{Estamos diciendo}: es PV. DE: V auxiliar
\forma{estar} + V auxiliado \forma{decir} en gerundio. 4. \forma{Estoy
aburrida}: no es PV. \forma{Dicen defenderme}: no es PV. 5. \forma{Iremos a
pagarle}: no es PV. \forma{Van a tener que entregar}:\footnote{En
§\ref{sec:11-cadenas} se describe el funcionamiento de las agrupaciones o
cadenas perifrásticas.} es PV. DE: V auxiliar\etq{1} \forma{ir} + \forma{a} + V
auxiliar\etq{2} \forma{tener} en infinitivo + \forma{que} + V auxiliado en
infinitivo. 6. \forma{Estamos pagando}: es PV. DE: V auxiliar \forma{estar} + V
auxiliado \forma{pagar} en gerundio. \forma{Vamos a pagar}: es PV. DE: V auxiliar
\forma{ir} + \forma{a} + V auxiliado en infinitivo.\pendiente{En 2,
\textit{iremos a comprar} figura como PV y, en 5, \textit{iremos a pagarle}, como
no PV: revisar}
\end{solucion}

\begin{solucion}{\ref{ej:11-vf}}
2. Verdadero. Los auxiliares de perífrasis solo pierden (parte de) su significado
cuando se combinan con una forma verbal no finita, pero en sus otros usos
continúan funcionando como verbos plenos. Ejemplos: \forma{Va todos los días a
la escuela} (V pleno con significado de movimiento); \forma{Va aprendiendo cada
día más} (V auxiliar que indica progreso del evento y no desplazamiento). 3.
Falso. Las perífrasis de participio solo aceptan la anteposición de los clíticos.
Ejemplo: \forma{El candidato ha rechazado el premio / lo ha rechazado /
*ha rechazádolo}. 4. Falso. Los auxiliares de perífrasis pueden construirse con
verbos impersonales. Ejemplo: \forma{Mañana puede llover todo el día}.
\end{solucion}

\begin{solucion}{\ref{ej:11-caracteristicas}}
2. No es PV porque se puede omitir el verbo principal en la interrogación. 3. Es
PV porque el verbo principal está reemplazado en la interrogación por el proverbo
\forma{hacer}. 4. Es PV porque es posible intercalar el adverbio
\forma{recientemente} entre los verbos que la componen. 5. No es PV porque
\forma{necesitar} mantiene su significado de V pleno, es decir que no está
desemantizado.
\end{solucion}

\begin{solucion}{\ref{ej:11-corpes}}
2. No se trata de una PV porque el participio \forma{descabellada} actúa como
adjetivo y funciona sintácticamente como el predicativo del V \forma{ser}. 3. Se
trata de una PV de \forma{ser} + participio que manifiesta la voz pasiva. 4. No
es una PV porque el gerundio funciona como un adverbio de modo y puede ser
reemplazado por \forma{así}. 5. \forma{Comenzó} y \forma{pude} no forman PV
porque son dos formas finitas; \forma{pude} forma PV con \forma{adivinar}. 6. Se
trata de una PV de \forma{comenzar} + \forma{a} + infinitivo que manifiesta el
inicio del evento.
\end{solucion}

\begin{solucion}{\ref{ej:11-pleno}}
2a. \sig{destruir}; 2b. \sig{comenzar abruptamente}. 3a. existencia; 3b.
anterioridad a otro evento del pasado. 4a. posesión de un atributo; 4b. voz
pasiva. 5a. ubicación; 5b. comienzo inminente.
\end{solucion}

\begin{solucion}{\ref{ej:11-significados}}
3. PV aspectual de cuantificación progresiva. 4. No es PV. 5. Es ambigua y posee
dos interpretaciones: i) no es PV porque indica dos eventos y admite conmutación
estructural: \forma{María va por el camino (mientras lee un libro) y aprende
gramática}; ii) PV aspectual de cuantificación progresiva. 6. PV aspectual de
fase continuativa. 7. PV de diátesis pasiva. 8. No es PV. 9. PV de modalidad
habitual. 10. PV aspectual de fase incoativa. 11. El contexto (la oposición entre
\forma{no volvió a rendir} y \forma{aprobó}) determina la interpretación de
\forma{volver} como PV aspectual de cuantificación reiterativa. 12. El contexto
(la oposición entre \forma{no vuelve a estudiar\ldots, sino para descansar})
determina la interpretación de \forma{volver} como V pleno de movimiento.
\end{solucion}

\begin{solucion}{\ref{ej:11-reformular}}
3. Mañana le voy a dar a mi compañero la contraseña. 4. Debe detenerse. 5. Aquí
suele nevar en invierno. 6. Ricardo se debe (de) haber dormido / puede haberse
dormido. 7. Enseguida me pongo / comienzo / empiezo a armar el equipo. 8. El
escritor acaba de terminar su última novela. 9. El escritor ha vivido durante
muchos años en San Juan, pero ahora vive en Santiago de Chile. 10. Mariela
volverá a ser mamá. 11. El profesor les pidió a sus alumnos que sigan / continúen
estudiando. 12. Por favor, ya dejen de pelear. 13. Todos sabemos que Miguel pudo
vencer a su rival. 14. Finalmente, Rosa pudo dormirse.
\end{solucion}

\begin{solucion}{\ref{ej:11-cadenas}}
2. \forma{Pudo} (posibilidad) + \forma{terminar de} (terminación) +
\forma{estudiar}. 3. \forma{Va a} (futuridad) + \forma{poder} (posibilidad) +
\forma{empezar a} (inicio) + \forma{estudiar}. 4. \forma{Deberían} (obligación) +
\forma{haber} (anterioridad) + \forma{empezado a} (inicio) + \forma{estudiar}. 5.
\forma{Tiene que} (obligación) + \forma{poder} (capacidad) + \forma{dejar de}
(interrupción / término) + \forma{comer}.
\end{solucion}

\begin{solucion}{\ref{ej:11-orden}}
\exref{ex:11-chomsky-b}: es gramatical porque la PV de modalidad precede a la de
voz pasiva. \exref{ex:11-chomsky-c}: es agramatical porque la PV de voz pasiva no
puede anteceder a ninguna otra en la oración. \exref{ex:11-chomsky-d}: es
gramatical porque la PV de tiempo compuesto precede a la aspectual reiterativa
con \forma{volver}. \exref{ex:11-chomsky-e}: es agramatical porque la PV
aspectual de fase continuativa no puede preceder a la de tiempo
compuesto.\pendiente{El original numera estas respuestas (70) a (73)}
\end{solucion}

\begin{solucion}{\ref{ej:11-tabla}}
\begin{center}
\resizebox{\linewidth}{!}{%
\begin{tabular}{@{}c *{6}{c} c@{}}
  \toprule
  & \multicolumn{6}{c}{Verbos auxiliares} & \\
  \cmidrule(lr){2-7}
  Ej. & \makecell{Mod. epistémica /\\habitualidad} & Temporalidad &
    \makecell{Mod.\\deóntica} & \makecell{Aspect.\\tipo II} &
    \makecell{Aspect.\\tipo I} & \makecell{Voz\\pasiva} & \makecell{Verbo\\auxiliado} \\
  \midrule
  2 & \textit{puede} & \textit{haber} & \textit{tenido que} & --- & \textit{dejar
    de} & \textit{ser} & \textit{atendido} \\
  3 & \textit{suelen} & --- & \textit{tener que} & \textit{volver a} & --- & ---
    & \textit{leer} \\
  4 & --- & \textit{han} & \textit{tenido que} & --- & --- & \textit{ser} &
    \textit{restaurados} \\
  5 & \textit{deben (de)} & \textit{haber} & --- & --- & --- & \textit{sido} &
    \textit{deterioradas} \\
  6 & --- & \textit{vamos a} & --- & --- & \textit{comenzar a} & --- &
    \textit{practicar} \\
  \bottomrule
\end{tabular}}
\end{center}
\end{solucion}

\begin{solucion}{\ref{ej:11-valores}}
2. Significado: modal deóntico (\forma{puede}) + aspectual continuativo
(\forma{seguir}) + de voz pasiva (\forma{siendo considerada}). La PV modal
indica que se admite o permite que Susana siga siendo considerada\ldots, por lo
que se considera deóntica. 3. Significado: modal epistémica (\forma{poder}) +
aspectual reiterativa (\forma{volver}) + pasiva (\forma{ser atendido}). La PV
modal se interpreta como \sig{tener la posibilidad de ser atendido\ldots}, por lo
que se considera epistémica.\pendiente{Según la jerarquía del capítulo,
\textit{poder} antepuesto a una aspectual con lectura de \textit{tener la
posibilidad} es deóntico: revisar} 4. Significado: témporo-aspectual
(\forma{ha}) + aspectual reiterativa (\forma{vuelto a}) + aspectual de
interrupción (\forma{dejar de tomar}).
\end{solucion}

\begin{solucion}{\ref{ej:11-entrar}}
2, 3, 5 y 6 son PPVV: no admiten conmutación estructural. La PV de
\forma{entrar} + infinitivo tiene valor aspectual de inicio. 4 no es PV: admite
conmutación estructural (\forma{entraron y robaron / entraron para robar}) y
\forma{entrar} actúa como verbo pleno.
\end{solucion}

\begin{solucion}{\ref{ej:11-locuciones}}
a) 2, 5 y 10 son PPVV y no aceptan conmutación estructural. Aceptan conmutación
léxica: 2. \forma{hemos asimilado}; 5. \forma{vamos a gozar}; 10. \forma{puede
existir}. 3, 4, 6, 7, 8 y 9 son LLVV: no aceptan ni conmutación léxica ni
conmutación estructural.\pendiente{El original numera «(1) hemos asimilado»} b)
Equivalencias: 3. \sig{extrañar}; 4. \sig{destrozar}; 6. \sig{informar}; 7.
\sig{considerar}; 8. \sig{reconocer}; 9. \sig{finalizar}.
\end{solucion}
